
\documentclass[a4paper]{article}

% Para usar vários arquivos com cabeçalho
\usepackage{standalone}

% Pacotes de formatação das fontes
\usepackage[brazil]{babel}
\usepackage[utf8]{inputenc}
\usepackage[T1]{fontenc}
\usepackage{fontspec}
\setmainfont[
    Path=fonts/,
    BoldItalicFont=*-bold-italic-2.ttf,
    BoldFont      =*-bold-2.ttf,
    ItalicFont    =*-italic.ttf]{calibri.ttf}


% Pacote para gerenciamento de espaçamentos
\usepackage[left=30mm,right=30mm,top=25mm,bottom=25mm]{geometry}
\usepackage{indentfirst}
\setlength{\parindent}{.5cm}
\setlength{\parskip}{1em}
\renewcommand{\baselinestretch}{1.5}

% Outros pacotes
\usepackage{xcolor}
\usepackage{float}
\usepackage{longtable}
\usepackage{amsmath,amssymb,amsthm}
\usepackage{multirow}
\usepackage{makecell}
\usepackage{enumitem}
\usepackage[labelfont=bf,small,font=sl,justification=centering,skip=2pt]{caption}
\captionsetup[table]{skip=2pt}


% Gerenciamento dos hyperlinks
\definecolor{dark-green}{HTML}{00781c}
\usepackage{hyperref}
\hypersetup{
  colorlinks=true,
  breaklinks=true,
  linkcolor=dark-green,
  citecolor=dark-green,
  urlcolor=dark-green,
  pdfsubject={Projeto Pedagógico BCC},
  pdftitle={Projeto Pedagógico BCC}
}
\urlstyle{same}

% Gerenciamento de figuras
\usepackage{graphicx}
\usepackage{tikz}

\usetikzlibrary{shapes,snakes,positioning,calc,arrows.meta,decorations.markings}
\tikzset{seta/.style={decoration={markings,mark=at position 1 with %
    {\arrow[scale=3,>=stealth]{>}}},postaction={decorate}}}

\newcommand{\disciplina}[6][yellow!20]{
\node [draw=black, fill=#1, very thick, rectangle, rounded corners, inner
sep=2pt, inner ysep=2pt, text width=#5pt, minimum height=40pt,
text centered, anchor=west] (#4) at (#2,#3) {
   \begin{minipage}{#5pt}
   \linespread{1.0}\selectfont
       \centering
       \scriptsize{#6}
   \end{minipage}
};
}
\newcommand{\disciplinab}[6][yellow!20]{
\node [draw=black, fill=#1, very thick, rectangle, rounded corners, inner
sep=2pt, inner ysep=2pt, text width=#5pt, minimum height=40pt,
text centered, anchor=west] (#4) at (#2,#3) {
   \begin{minipage}{#5pt}
   \linespread{1.0}\selectfont
       \centering
       \small{#6}
   \end{minipage}
};
}

% Espaçamentos nas listas
\setlist[itemize]{noitemsep,topsep=0pt}
\setlist[enumerate]{noitemsep,topsep=0pt}

% Paleta de cores daltônico-friendly
\definecolor{nred}{rgb}{0.88, 0.28, 0.33}
\definecolor{nblue}{rgb}{0.34, 0.74, 0.96}
\definecolor{nyellow}{rgb}{0.8, 0.86, 0.38}
\definecolor{npurple}{rgb}{0.25, 0.34, 0.93}
\definecolor{ngreen}{rgb}{0.45, 0.95, 0.66}


% Outros
\usepackage{sectsty}
\sectionfont{\MakeUppercase}
\subsectionfont{\MakeUppercase} 

\usepackage{tocloft}
\renewcommand\cftsecleader{\cftdotfill{\cftdotsep}}

\newcommand\T{\rule{0pt}{5.6ex}}       % Top strut
\newcommand\B{\rule[-1.2ex]{0pt}{0pt}} % Bottom strut

\newcommand{\red}[1]{\textcolor{red}{#1}}


% Gerenciamento da bibliografia
\usepackage[style=abnt-numeric,sorting=none,maxbibnames=99]{biblatex}
\setlength\bibitemsep{0pt}
\defbibheading{bibliography}[\bibname]{}
\defbibnote{space}{\vskip -0.5cm}

\let\OLDprintbibliography\printbibliography
\renewcommand\printbibliography{
    \OLDprintbibliography[prenote=space]
}

\newcommand{\incluircurso}[1]{\input{#1}\cleardoublepage}

\addbibresource{obrigatorias/MCTA001.bib}
\addbibresource{obrigatorias/MCTA001-compl.bib}
\addbibresource{obrigatorias/MCTA002.bib}
\addbibresource{obrigatorias/MCTA002-compl.bib}
\addbibresource{obrigatorias/MCTA003.bib}
\addbibresource{obrigatorias/MCTA003-compl.bib} 
\addbibresource{obrigatorias/MCTA004.bib}
\addbibresource{obrigatorias/MCTA004-compl.bib}
\addbibresource{obrigatorias/MCTA006.bib}
\addbibresource{obrigatorias/MCTA006-compl.bib}
\addbibresource{obrigatorias/MCTA007.bib}
\addbibresource{obrigatorias/MCTA007-compl.bib}
\addbibresource{obrigatorias/MCTA008.bib}
\addbibresource{obrigatorias/MCTA008-compl.bib}
\addbibresource{obrigatorias/MCTA009.bib}
\addbibresource{obrigatorias/MCTA009-compl.bib}
\addbibresource{obrigatorias/MCTA014.bib}
\addbibresource{obrigatorias/MCTA014-compl.bib}
\addbibresource{obrigatorias/MCTA015.bib}
\addbibresource{obrigatorias/MCTA015-compl.bib}
\addbibresource{obrigatorias/MCTA016.bib}
\addbibresource{obrigatorias/MCTA016-compl.bib}
\addbibresource{obrigatorias/MCTA017.bib}
\addbibresource{obrigatorias/MCTA017-compl.bib}
\addbibresource{obrigatorias/MCTA018.bib}
\addbibresource{obrigatorias/MCTA018-compl.bib}
\addbibresource{obrigatorias/MCTA022.bib}
\addbibresource{obrigatorias/MCTA022-compl.bib}
\addbibresource{obrigatorias/MCTA023.bib}
\addbibresource{obrigatorias/MCTA023-compl.bib}
\addbibresource{obrigatorias/MCTA024.bib}
\addbibresource{obrigatorias/MCTA024-compl.bib}
\addbibresource{obrigatorias/MCTA025.bib}
\addbibresource{obrigatorias/MCTA025-compl.bib}
\addbibresource{obrigatorias/MCTA026.bib}
\addbibresource{obrigatorias/MCTA026-compl.bib}
\addbibresource{obrigatorias/MCTA028.bib}
\addbibresource{obrigatorias/MCTA028-compl.bib}
\addbibresource{obrigatorias/MCTA033.bib}
\addbibresource{obrigatorias/MCTA033-compl.bib}
\addbibresource{obrigatorias/MCTA037.bib}
\addbibresource{obrigatorias/MCTA037-compl.bib}
\addbibresource{obrigatorias/MCTAnova1.bib}
\addbibresource{obrigatorias/MCTAnova1-compl.bib}
\addbibresource{obrigatorias/MCTAnova3.bib}
\addbibresource{obrigatorias/MCTAnova3-compl.bib}
\addbibresource{obrigatorias/MCTAnova4.bib}
\addbibresource{obrigatorias/MCTAnova4-compl.bib}
\addbibresource{obrigatorias/MCTAnova5.bib}
\addbibresource{obrigatorias/MCTAnova5-compl.bib}


\begin{document}


\begin{tabular}{rp{.5\textwidth}r}
    \multirow{2}{*}{\includegraphics[scale=.03]{imagens/brasao.png}} & 
    \centering \Large{Ministério da Educação} & \multirow{2}{*}{\includegraphics[scale=2.3]{imagens/ufabc-logo.png}} \\
    & \centering \Large{Universidade Federal do ABC} &\\
\end{tabular}

\vspace{9cm}

\begin{center}
    \Large{PROJETO PEDAGÓGICO DO CURSO DE \\\textbf{BACHARELADO EM CIÊNCIA DA COMPUTAÇÃO}}
\end{center}

\vspace{9cm}

\begin{center}
    SANTO ANDRÉ \\ 2023
\end{center}

\newpage

\begin{tabular}{l}
    \textbf{Reitor da UFABC}\\
    Prof.\ Dr.\ Dácio Roberto Matheus\\
    \\
    \textbf{Pró-Reitoria de Graduação}\\
    Profa.\ Dra.\ Fernanda Graziella Cardoso -- Pró-reitora\\
    Prof.\ Dr.\ Marcelo Salvador Caetano -- Pró-reitor adjunto\\
    \\
    \textbf{Diretoria do Centro de Matemática, Computação e Cognição}\\
    Profa.\ Dra.\ Tatiana Lima Ferreira -- Diretora\\
    Prof.\ Dr.\ Maurício Richartz -- Vice-diretor\\
    \\
    \textbf{Coordenação do Curso de Bacharelado em Ciência da Computação}\\
    Prof.\ Dr.\ Jerônimo Cordoni Pellegrini -- Coordenador\\
    Prof.\ Dr.\ Gustavo Sousa Pavani -- Vice-coordenador\\
    \\
    \textbf{Equipe de Trabalho}\\
    Prof.\ Dr.\ Alexandre Donizeti Alves\\
    Prof.\ Dr.\ Aritanan Borges Garcia Gruber\\
    Profa.\ Dra.\ Carla Lopes Rodriguez\\
    Profa.\ Dra.\ Carla Negri Lintzmayer\\
    Prof.\ Dr.\ Carlos da Silva dos Santos\\
    Profa.\ Dra.\ Denise Hideko Goya\\
    Prof.\ Dr.\ Emílio de Camargo Francesquini\\
    Prof.\ Dr.\ Francisco de Assis Zampirolli\\
    Prof.\ Dr.\ Fabrício Olivetti de França\\
    Prof.\ Dr.\ Gustavo Sousa Pavani\\
    Prof.\ Dr.\ Jerônimo Cordoni Pellegrini\\
    Prof.\ Dr.\ João Marcelo Borovina Josko\\
    Prof.\ Dr.\ Marcio Katsumi Oikawa\\
    Prof.\ Dr.\ Maycon Sambinelli\\
    Prof.\ Dr.\ Vladimir Emiliano Moreira Rocha\\
\end{tabular}



\newpage

\tableofcontents

\newpage

\section{Dados da Instituição}

\textbf{Nome da Unidade}: Fundação Universidade Federal do ABC

\textbf{CNPJ}: 07.722.779/0001-06

\textbf{Lei de Criação}: Lei nº 11.145, de 26 de julho de 2005, publicada no
D.O.U. em 27 de julho de 2005, alterada pela Lei nº 13.110, de 25 de março de
2015, publicada no D.O.U. em 26 de março de 2015\footnote{Disponível em:
\url{http://www.ufabc.edu.br/a-ufabc/documentos/lei-de-criacao-da-ufabc}.
Acesso em 12 de dezembro de 2022.}.


\newpage
\section{Dados do Curso}
\label{sec:dados_curso}

\textbf{Curso}: Bacharelado em Ciência da Computação.

\textbf{Diplomação}: Bacharel em Ciência da Computação.

\textbf{Carga horária total do curso}: \textcolor{red}{TODO: corrigir aqui 3.200 horas}.

\textbf{Tempo esperado de integralização}: cinco anos.  

\textbf{Tempo máximo de integralização}: oito anos. Devido a características
específicas da UFABC, o cálculo do tempo mínimo depende do bacharelado
interdisciplinar de acesso. Ver maiores detalhes na resolução
correspondente\footnote{Disponível em:
\url{https://www.ufabc.edu.br/images/consepe/resolucoes/resolucao_166_-_desligamento.pdf}.
Acesso em 12 de dezembro de 2022.}.

\textbf{Estágio}: o estágio curricular não é obrigatório.

\textbf{Trabalho de conclusão de curso}: obrigatório no último ano do curso.

\textbf{Turnos de oferta}: matutino e noturno.

\textbf{Número de vagas por turno}: 70 (total de 140 vagas anuais).

\textbf{Câmpus de oferta}: Câmpus Santo André

\textbf{Endereço}: Av. dos Estados, 5.001 - Bairro Santa Terezinha. Santo André
- SP. CEP 09.210-580.

\textbf{Atos legais}: Lei de criação da UFABC. Resolução ConsEPE nº 195 de 2015
que aprova a revisão do projeto pedagógico. Curso reconhecido pela Portaria MEC
nº 406, de 11 de outubro de 2011, publicada no DOU em 14 de outubro de 2011.


\newpage
\section{Apresentação}
\label{sec:apresentacao}

A Fundação Universidade Federal do ABC (UFABC) é uma fundação pública criada
pela Lei nº 11.145 de 26 de julho de 2005, sancionada pelo Presidente da
República e publicada no Diário Oficial da União (DOU) em 27 de julho de 2005 e
alterada pela Lei nº 13.110, de 25 de março de 2015, publicada no DOU em 26 de
março de 2015.
É uma instituição de ensino superior, extensão e pesquisa, vinculada ao
Ministério da Educação (MEC), com sede e foro no Município de Santo André,
situada na Avenida dos Estados, 5001, bairro Santa Terezinha, Santo André, CEP
09210-580, no Estado de São Paulo e com limite territorial de atuação
multicampi na região do ABC paulista, nos termos do Artigo 2º da lei
mencionada.

A UFABC possui autonomia administrativa, didático-científica, gestão financeira
e disciplinar, rege-se pela legislação federal que lhe é pertinente, pelo
Regimento dos Órgãos da Administração Superior e das Unidades Universitárias e
pelas Resoluções de seus Órgãos.

A instituição foi criada para atender a um anseio antigo da região do ABC
paulista por uma universidade pública e de qualidade.
A UFABC busca ser reconhecida como uma referência no panorama nacional e
internacional, por meio de sua atenção às demandas regionais, produzindo
pesquisas e formando profissionais de alta qualidade para enfrentá-las.
Sua missão é \textit{``facilitar e induzir a interdisciplinaridade, promovendo
a visão sistêmica e a apropriação do conhecimento pela sociedade, sem
esmorecimento da rigorosa cultura disciplinar''}.
Para esse propósito, a UFABC procura ter um olhar voltado para o mundo e, ao
mesmo tempo, procura caminhar lado a lado com a sociedade e o setor produtivo.

Nesse propósito, a atuação acadêmica da UFABC se dá através de cursos de
graduação, pós-graduação e extensão, visando a formação e o aperfeiçoamento de
recursos humanos solicitados para o progresso da sociedade brasileira.
Além disso, a instituição promove e estimula a pesquisa científica, tecnológica
e a produção de pensamento original no campo da ciência e da tecnologia.

Todos os alunos de graduação da UFABC ingressam por meio de algum curso
interdisciplinar:
\begin{itemize}
    \item Bacharelado em Ciência e Tecnologia (BC\&T);
    \item Bacharelado em Ciências e Humanidades (BC\&H);
    \item Licenciatura em Ciências Naturais Exatas (LCNE); ou
    \item Licenciatura em Ciências Humanas (LCH).
\end{itemize}
Além desses, são ofertados 20 cursos de bacharelado específicos e 5 cursos de
licenciatura específicas.



\subsection{O curso de Bacharelado em Ciência da Computação da UFABC}

O curso de Bacharelado em Ciência da Computação (BCC), previsto no Projeto
Pedagógico Institucional da UFABC (PPI-UFABC), faz parte do planejamento global
da instituição, que tem entre seus objetivos tornar-se um pólo produtor de
conhecimento de nível nacional e internacional no âmbito das ciências, cultura
e artes.

O BCC está sediado no câmpus Santo André e iniciou seu funcionamento a partir
do Edital de Vestibular ocorrido em 02 de maio de 2006, publicado no DOU,
Seção 3, nº 83, 3 de maio de 2006, pg.\ 25.

O BCC tem a duração mínima de cinco anos, podendo ser concluído em prazo menor
a depender do desempenho do aluno e do regime de matrículas da UFABC.
A duração máxima do curso é de oito anos, conforme a Resolução ConsEP nº 166,
de 8 de outubro de 2013.
Deve-se atentar ao prazo máximo de 18 quadrimestres para integralização do
BC\&T, conforme Resolução ConsEPE nº 166, de 8 de outubro de 2013.

A admissão no BCC pode ser realizada por discentes que estão cursando ou já
concluíram o BC\&T.
As disciplinas sugeridas na matriz curricular do BC\&T podem ser cursadas
paralelamente às disciplinas sugeridas na matriz curricular do BCC.
Apesar disso, a colação de grau no BCC está vinculada à colação de grau no
BC\&T, de modo que o aluno que desejar colar grau no BCC já deve possuir o grau
de Bacharel em Ciência e Tecnologia.
A colação de grau de ambos os cursos também pode ser realizada de forma conjunta.

Além de garantir aos egressos uma sólida e abrangente formação em Ciência da
Computação por meio de suas disciplinas obrigatórias e de opção limitada, o
curso se compromete com atividades complementares à sua formação, tais como
monitoria acadêmica, iniciação científica e atividades extensionistas.

Este projeto pedagógico baseia seu conteúdo na integração dos seguintes
documentos reguladores:
\begin{itemize}
    \item PPI-UFABC de 2017;
    \item Projeto pedagógico do BC\&T (2023);
    \item Diretrizes curriculares nacionais para os cursos de graduação da Computação (2016);
    \item Referenciais de formação para cursos de Graduação em Computação (2017);
    \item Plano de Desenvolvimento Institucional (PDI) da UFABC (2013-2022).
\end{itemize}

No terceiro quadrimestre de 2010, formou-se a primeira turma do BCC e, em março
de 2011, a comissão designada pelo INEP/MEC emitiu parecer favorável ao
reconhecimento do curso, atribuindo ao mesmo o conceito máximo 5 (cinco).
Na aplicação do Exame Nacional de Desempenho de Estudantes (Enade) realizado em
2021, o curso também obteve conceito máximo 5 (cinco).


\newpage
\section{Perfil do Curso}
\label{sec:perfil_curso}

O BCC da UFABC propõe formar profissionais com carácter interdisciplinar e
multidisciplinar, com formação teórica consistente e vivência prática que
permita contribuir para o desenvolvimento científico e tecnológico da Ciência
da Computação, atuando profissionalmente em empresas de tecnologia, em pesquisa
científica ou em ações empreendedoras.

Além de uma formação básica sólida e uma proposta de desenvolvimento ético e
científico, o curso promove fortemente uma construção interdisciplinar, em
consonância com PPI-UFABC.
Os egressos do curso podem atuar em nível regional, nacional e internacional,
atendendo à crescente demanda por profissionais qualificados nas diversas áreas
em que a Ciência da Computação pode atuar.

A Computação está presente na rotina da população em praticamente todas as suas
atividades sociais, econômicas e científicas.
Podemos facilmente identificar a influência de algoritmos e recursos
computacionais em diversas atividades comuns, tais como ler notícias,
comunicar-se com outras pessoas, viajar, trabalhar, estudar, etc.
Dispositivos computacionais estão presentes em eletrodomésticos, veículos,
telefones celulares, televisores e computadores, entre outros.
A Ciência da Computação é certamente uma das áreas de futuro mais promissor,
abrindo várias oportunidades de desenvolvimento tecnológico e alimentando
iniciativas empreendedoras que buscam soluções para problemas gerais e
específicos da sociedade.
A demanda por profissionais é reconhecidamente alta e com tendência de
expansão, necessitando de cursos de formação que contribuam para atender de
forma qualificada a essa perspectiva de crescimento.

A estrutura curricular do BCC se baseia em vários documentos de referência:
\begin{itemize}
    \item Diretrizes Curriculares Nacionais (DCN) dos cursos da área de
    Computação;
    \item Proposta curricular das associações:
    \begin{itemize}
        \item ACM (\textit{Association for Computing Machinery});
        \item IEEE-CS (\textit{IEEE Computer Society});
        \item SBC (Sociedade Brasileira de Computação).
    \end{itemize}
\end{itemize}

O BC\&T, curso de ingresso de todo aluno do BCC, contribui com a formação
básica e divide-se em seis eixos didáticopedagógicos estruturantes:
\begin{itemize}
    \item Estrutura da Matéria;
    \item Energia;
    \item Processos de Transformação;
    \item Representação e Simulação;
    \item Informação e Comunicação;
    \item Humanidades.
\end{itemize}
Assim, aliada ao BC\&T, a estrutura curricular do BCC abrange diversas áreas de
formação.


\subsection{Justificativa de oferta do curso}

A UFABC está localizada na região conhecida como ABC Paulista, apelido que faz
referência às cidades de Santo \textbf{A}ndré, São \textbf{B}ernardo do Campo e
São \textbf{C}aetano do Sul, e também faz parte da região metropolitana de São
Paulo (RMSP).
A RMSP é altamente urbanizada (98\%) formada por 39 municípios e uma população
próxima de 22 milhões de habitantes (2021), que a faz figurar entre as dez mais
populosas do mundo\footnote{Dados disponíveis em
\url{https://perfil.seade.gov.br/}. Acesso em 14 de dezembro de 2022.}. 

Do ponto de vista econômico, a RMSP é considerada o maior pólo de riqueza do
Brasil, com PIB per capita no valor de R\$ 56.649,03 (2018).
A atividade econômica está fortemente ligada à prestação de serviços (85,5\%),
embora o setor industrial também tenha relevância (14,3\%), sendo grande a
contribuição do ABC Paulista. 
Do ponto de vista educacional, é uma região em que mais da metade (57,5\%) da
população jovem entre 18 e 24 anos possui, no mínimo, o Ensino Médio completo
(censo 2010).
É também uma região com grande número de escolas e faculdades, públicas e
privadas. 

A Computação é uma das áreas mais promissoras em termos crescimento e
desenvolvimento.
Praticamente todos os setores utilizam recursos computacionais para automatizar
tarefas, desenvolver produtos, otimizar a utilização e monitoramento de
recursos, inovar, planejar políticas de expansão, controlar atividades, etc.
Durante a pandemia de COVID-19, foi uma das poucas áreas que apresentou
crescimento e permitiu que muitas atividades econômicas e sociais pudessem ser
preservadas, apesar das dificuldades e restrições sanitárias.
Segundo levantamento da Associação Brasileira das Empresas de Tecnologia da
Informação e Comunicação (Brasscom) realizado em 2021, a demanda não atendida
por profissionais no Brasil deve atingir 420 mil vagas até 2024, sendo que
formam-se aproximadamente 46 mil por ano.

Nesse contexto, o ABC Paulista pode ser visto como uma região estratégica para
o apoio ao desenvolvimento tecnológico local e nacional.
O ABC é uma região com forte participação industrial, conurbada em uma área com
forte demanda por serviços.
Além disso, é uma região com alto índice educacional, integrada à RMSP e ao
Brasil por meio de grandes rodovias, grandes aeroportos, ferrovias e o porto de
Santos, o maior da América Latina.
É uma região estratégica para implantação de empresas nacionais e
internacionais, de diversos setores sociais e econômicos.

A Computação é uma das áreas de conhecimento mais presente e influente na vida
de empresas e pessoas.
Encontramos técnicas, teorias, produtos e metodologias associadas à Ciência da
Computação em diversas iniciativas empresariais, políticas, sociais e
tecnológicas.
A busca por profissionais qualificados é uma necessidade de diversas entidades
que buscam inovação, otimização de recursos, pesquisa e desenvolvimento.

Outra característica da Ciência da Computação é sua aplicabilidade, capaz de
contribuir com diversas áreas de conhecimento, o que lhe garante alta
capacidade interdisciplinar e integradora.
A implantação do BCC, sob essa visão, é naturalmente identificada aos
princípios norteadores da UFABC e às necessidades das comunidades local,
regional e nacional.



\newpage
\section {Objetivos do Curso}

\subsection{Objetivo Geral}
Formação de profissionais com sólido conhecimento científico e tecnológico na área de Computação.

\subsection{Objetivos Específicos}
\begin{itemize}
	\item Incentivar o perfil pesquisador do estudante, visando promover o desenvolvimento científico e tecnológico da Ciência da Computação;
	\item Preparar o estudante para atuar profissionalmente em organizações, com espírito empreendedor e com responsabilidade social;
	\item Proporcionar atividades acadêmicas que estimulem a interdisciplinaridade, bem como a aplicação e renovação dos conhecimentos e habilidades de forma independente e inovadora, nos diversos contextos da atuação profissional;
	\item Formar estudantes que possam estar em sintonia com a nova realidade e necessidade do aprendizado contínuo e autônomo, exigido pela sociedade do conhecimento e organizações dos dias atuais;
	\item Promover no estudante uma postura ética e socialmente comprometida de seu papel e de sua contribuição no avanço científico, tecnológico e social do País.
	
\end{itemize}
Com base nesses objetivos, pode-se definir que o bacharel em Ciência da Computação da
UFABC deverá conhecer os fundamentos de sua ciência, suas raízes históricas e suas
interligações com outras ciências.



\newpage
\section {Requisito de Acesso}

\subsection{Formas de Acesso ao Curso}
O processo seletivo para acesso aos Cursos de Graduação da Universidade Federal do ABC é anual, e inicialmente realizado pelo Sistema de Seleção Unificado (SISU), do MEC, onde as vagas oferecidas serão preenchidas em uma única fase, baseado no resultado do Exame Nacional do Ensino Médio (ENEM).

O ingresso nos cursos de formação específica, após a conclusão dos bacharelados
interdisciplinares, se dá por seleção interna, segundo a Resolução ConsEPE, no. 31, de 4 de agosto de 2009. Sendo assim, o ingresso ao Bacharelado em Ciência da Computação é realizado após o ingreso no Bacharelado em Ciência e Tecnologia.

Existe ainda a possibilidade de transferência, facultativa ou obrigatória, de alunos de outras Instituições de Ensino (IES) para o BCC. No primeiro caso, mediante transferência de alunos de cursos afins, quando da disponibilidade de vagas, através de processo seletivo interno (art. 49 da Lei n° 9.394, de 1996 e Resolução ConsEPE no. 174 de 24 de abril de 2014); para o segundo, por \textit{transferências ex officio} previstas em normas específicas (art. 99 da Lei 8.112 de 1990, art. 49 da Lei 9.394 de 1996 regulamentada pela Lei 9.536 de 1997 e Resolução ConsEPE no. 10 de 2008)

\subsection{Regime de Matrícula}
O processo de matrículas em disciplinas é conduzido de forma unificada pela Pró-Reitoria de Graduação (Prograd) da UFABC. Antes do início de cada quadrimestre letivo, cada aluno(a) deve solicitar a sua matrícula, indicando as disciplinas que deseja cursar no quadrimestre correspondente. O período de matrícula é determinado pelo calendário da UFABC definido anualmente pela Comissão de Graduação.

A matrícula de alunos ingressantes é realizada de forma automática e obrigatória, obedecendo à matriz curricular do bacharelado interdisciplinar de ingresso. A partir do quadrimestre letivo seguinte, o(a) aluno(a) entra no regime de matrícula regular, obedecendo ao procedimento citado anteriormente. Alunos ingressantes devem cursar, obrigatoriamente, o mínimo de nove créditos no quadrimestre de ingresso. 

Por não apresentarem pré-requisitos, todas as disciplinas podem ser solicitadas livremente e a qualquer momento no processo de matrícula. Apesar disso, deve-se ressaltar que cada disciplina possui uma lista de recomendações, que expõe disciplinas que desejavelmente deveriam ter sido cursadas anteriormente. Embora não exista o bloqueio formal do pré-requisito, é importante que cada estudante considere a lista de recomendações como um elemento orientador que busca auxiliar o cumprimento bem-sucedido da matriz curricular.

É essencial ressaltar que o número de vagas e turmas é limitado, e o preenchimento de vagas na matrícula segue os critérios de seleção adotados pela Prograd. Em casos particulares (como disciplinas de Trabalho de Conclusão de Curso, Estágios ou outras), os pedidos de matrícula são ainda analisados pela coordenação do BCC, que poderá autorizá-los, ou não, dentro de seus critérios de adequação e viabilidade pedagógica.

É importante ainda que o(a) estudante observe os critérios de permanência do curso e jubilação (desligamento), regulados pela Resolução ConsEPE no. 166, de 08 de outubro de 2013.


\newpage
\section{Perfil do egresso}
\label{sec:perfil_egresso}

O BCC baseia-se em dois conjuntos fundamentais de documentos para a composição de sua proposta pedagógica e curricular e formação do perfil de egresso:
\begin{itemize}
	\item Diretrizes curriculares nacionais (DCNs) para cursos na área de Computação (ver Seção \ref{subsec:fund_legal});
	\item Projeto pedagógico institucional da UFABC (PPI).
\end{itemize}

Em relação às DCNs da área de Computação, a estrutura curricular se orienta nas exigências quanto ao "Perfil Geral de Egressos de Cursos na Área de Computação" e o "Perfil Específico de Egressos de Cursos de Bacharelado em Ciência da Computação", listados a seguir. 

O BCC trabalha a formação de seus egressos para que:
\begin{itemize}
	\item possuam sólida formação em Ciência da Computação e Matemática que os capacitem a construir aplicativos de propósito geral, ferramentas e infraestrutura de software de sistemas de computação e de sistemas embarcados, gerando conhecimento científico e inovação;
	\item desenvolvam visão global e interdisciplinar de sistemas e entendam que esta visão transcende os detalhes de implementação dos vários componentes e os conhecimentos	dos domínios de aplicação;
	\item conheçam a estrutura dos sistemas de computação e os processos envolvidos na sua construção e análise;
	\item dominem os fundamentos teóricos da área de Computação e como eles influenciam a prática profissional;
	\item sejam capazes de agir de forma reflexiva na construção de sistemas de computação, compreendendo o seu impacto direto ou indireto sobre as pessoas e a sociedade;
	\item sejam capazes de criar soluções, individualmente ou em equipe, para problemas complexos caracterizados por relações entre domínios de conhecimento e de aplicação;
	\item reconheçam o caráter fundamental da inovação e da criatividade e compreendam as perspectivas de negócios e oportunidades relevantes.
\end{itemize}

    Além disso e de forma mais ampla, o BCC trabalha seus egressos para desenvolver a capacidade de: 
\begin{itemize}
	\item identificar problemas que tenham solução algorítmica;
	\item conhecer os limites da computação;
	\item resolver problemas usando ambientes de programação;
	\item tomar decisões e inovar, com base no conhecimento do funcionamento e
	\item características técnicas de hardware e da infraestrutura de software dos sistemas de computação consciente dos aspectos éticos, legais e dos impactos ambientais decorrentes;
	\item compreender e explicar as dimensões quantitativas de um problema;
	\item gerir a sua própria aprendizagem e desenvolvimento, incluindo a gestão de tempo e competências organizacionais;
	\item preparar e apresentar seus trabalhos e problemas técnicos e suas soluções para audiências diversas, em formatos apropriados (oral e escrito);
	\item avaliar criticamente projetos de sistemas de computação;
	\item adequar-se rapidamente às mudanças tecnológicas e aos novos ambientes de trabalho;
	\item ler textos técnicos na língua inglesa;
	\item empreender e exercer liderança, coordenação e supervisão na sua área de atuação profissional;
	\item realizar trabalho cooperativo e entender os benefícios que este pode produzir.
	
\end{itemize}


De forma complementar, o BCC trabalha sua estrutura curricular e suas ações em conformidade com as três principais políticas institucionais previstas no PPI da UFABC: interdisciplinaridade, excelência e inclusão social. Por meio das iniciativas institucionais da UFABC, o aluno possui ampla liberdade para complementar sua formação em diversas áreas de conhecimento trabalhadas pelos cursos da universidade. Com suas várias oportunidades de integração com iniciativas de pesquisa e extensão, os alunos também encontram oportunidades para participar de projetos de pesquisa científica e tecnológica, além de acesso a equipamentos e técnicas avançadas de pesquisa moderna. Utilizando diversas iniciativas de inclusão social, a UFABC, dentro de suas possibilidades orçamentárias e legais, busca ampliar seu alcance para a comunidade local, promovendo ações que buscam democratizar o acesso ao ensino superior, compartilhar os resultados de iniciativas científicas, e abrir oportunidades para alunos em situação de vulnerabilidade.

Com isso e em consonância com as DCNs da área de Computação, também é papel do curso garantir que seu egresso seja dotado:

\begin{itemize}
	\item de conhecimento das questões sociais, profissionais, legais, éticas, políticas e humanísticas;
	\item da compreensão do impacto da computação e suas tecnologias na sociedade no que concerne ao atendimento e à antecipação estratégica das necessidades da sociedade;
	\item de visão crítica e criativa na identificação e resolução de problemas contribuindo para o desenvolvimento de sua área;
	\item da capacidade de atuar de forma empreendedora, abrangente e cooperativa no atendimento às demandas sociais da região onde atua, do Brasil e do mundo;
	\item de utilizar racionalmente os recursos disponíveis de forma transdisciplinar;
	\item da compreensão das necessidades da contínua atualização e aprimoramento de suas competências e habilidades;
	\item da capacidade de reconhecer a importância do pensamento computacional na vida cotidiana, como também sua aplicação em outros domínios e ser capaz de aplicá-lo em circunstâncias apropriadas; e
	\item da capacidade de atuar em um mundo de trabalho globalizado.
\end{itemize}
        
O profissional egresso do curso de Bacharelado em Ciência da Computação da UFABC poderá atuar tanto na análise e desenvolvimento de sistemas, em software ou hardware, como na pesquisa cientírica em Ciência da Computação.



\newpage
\section{Organização curricular}
\label{sec:organizacao_curricular}

\subsection{Fundamentação legal}
\label{subsec:fund_legal}

A seguir são elencados os documentos legais externos (Diretrizes
Curriculares Nacionais, Leis, Decretos, Resoluções, Pareceres,
Portarias, Normativas etc.), de ordem federal, bem como os internos
(Projeto Pedagógico, Plano de Desenvolvimento Institucional) que
fundamentam a estrutura curricular do curso de Bacharelado em Ciência
da Computação da UFABC:

\begin{itemize}
    
    \item BRASIL. Presidência da República. Casa Civil. Subchefia para Assuntos
    Jurídicos. Lei 9.394, de 20 de dezembro de 1996. Estabelece as diretrizes e
    bases da educação nacional.
    Disponível em: \url{https://www.planalto.gov.br/ccivil_03/Leis/L9394.htm}. 
    Acesso em: 20 de janeiro de 2023.
    
    \item BRASIL. Presidência da República. Casa Civil. Subchefia para Assuntos
    Jurídicos. Lei 10.639, de 9 de janeiro de 2003. Altera a Lei 9.394, de 20
    de dezembro de 1996, que estabelece as diretrizes e bases da educação
    nacional, para incluir no currículo oficial da Rede de Ensino a
    obrigatoriedade da temática ``História e Cultura Afro-Brasileira'', e dá
    outras providências. 
    Disponível em: \url{http://www.planalto.gov.br/ccivil_03/leis/2003/l10.639.htm}.
    Acesso em: 20 de janeiro de 2023.
    
    \item BRASIL. Presidência da República. Casa Civil. Subchefia para Assuntos
    Jurídicos. Lei 11.645, de 10 de março de 2008. Altera a Lei 9.394, de 20 de
    dezembro de 1996, modificada pela Lei 10.639, de 9 de janeiro de 2003, que
    estabelece as diretrizes e bases da educação nacional, para incluir no
    currículo oficial da rede de ensino a obrigatoriedade da temática
    ``História e Cultura Afro-Brasileira e Indígena''. 
    Disponível em: \url{http://www.planalto.gov.br/ccivil_03/_ato2007-2010/2008/lei/l11645.htm}.
    Acesso em: 20 de janeiro de 2023.
    
    \item BRASIL. Ministério da Educação. Conselho Nacional de Educação.
    Conselho Pleno. Resolução 1, de 17 de junho de 2004. Institui Diretrizes
    Curriculares Nacionais para a Educação das Relações Étnico Raciais e para o
    Ensino de História e Cultura Afro-Brasileira e Africana. 
    Disponível em: \url{http://portal.mec.gov.br/cne/arquivos/pdf/res012004.pdf}. 
    Acesso em: 20 de janeiro de 2023.
    
    \item BRASIL. Ministério da Educação. Secretaria da Educação Superior.
    Referenciais Orientadores para os Bacharelados Interdisciplinares e
    Similares. 2010.
    Disponível em: \url{http://www.ufabc.edu.br/images/stories/comunicacao/bacharelados-interdisciplinares_referenciais-orientadores-novembro_2010-brasilia.pdf}.
    Acesso em: 20 de janeiro de 2023.
    
    \item BRASIL. Ministério da Educação. Conselho Nacional de Educação.
    Referenciais orientadores para os Bacharelados Interdisciplinares e
    Similares das Universidades Federais. Câmara de Educação Superior. Parecer
    CNE/CES 266, de 5 de julho de 2011.
    Disponível em:
    \url{http://portal.mec.gov.br/index.php?option=com_content&view=article&id=16418&Itemid=866}
    Acesso em: 20 de janeiro de 2023.
    
    \item BRASIL. Presidência da República. Casa Civil. Subchefia para Assuntos
    Jurídicos. Decreto 5.626, de 22 de dezembro de 2005. Regulamenta a Lei
    10.436, de 24 de abril de 2002, que dispõe sobre a Língua Brasileira de
    Sinais - Libras, e o art. 18 da Lei 10.098, de 19 de dezembro de 2000.
    Disponível em: \url{https://www.planalto.gov.br/ccivil_03/_Ato2004-2006/2005/Decreto/D5626.htm}.
    Acesso em: 20 de janeiro de 2023.
    
    \item  BRASIL. Presidência da República. Casa Civil. Subchefia para
    Assuntos Jurídicos. Lei 9.795, de 27 de abril de 1999. Dispõe sobre a
    educação ambiental, institui a Política Nacional de Educação Ambiental e dá
    outras providências. 
    Disponível em: \url{http://www.planalto.gov.br/ccivil_03/leis/l9795.htm}.
    Acesso em: 20 de janeiro de 2023.
    
    \item BRASIL. Presidência da República. Casa Civil. Subchefia para Assuntos
    Jurídicos. Decreto 4.281, de 25 de junho de 2002. Regulamenta a Lei 9.795,
    de 27 de abril de 1999, que institui a Política Nacional de Educação
    Ambiental, e dá outras providências. 
    Disponível em: \url{http://www.planalto.gov.br/ccivil_03/decreto/2002/D4281.htm}.
    Acesso em: 20 de janeiro de 2023.
    
    \item BRASIL. Ministério da Educação. Conselho Nacional de Educação. Câmara
    de Educação Superior. Resolução 2, de 18 de junho de 2007. Dispõe sobre
    carga horária mínima e procedimentos relativos à integralização e duração
    dos cursos de graduação, bacharelados, na modalidade presencial. 
    Disponível em: \url{http://portal.mec.gov.br/cne/arquivos/pdf/2007/rces002_07.pdf}.
    Acesso em: 20 de janeiro de 2023.
    
    \item BRASIL. Ministério da Educação. Gabinete do Ministro. Portaria
    Normativa 40, de 12 de dezembro de 2007. Institui o e-MEC, sistema
    eletrônico de fluxo de trabalho e gerenciamento de informações relativas
    aos processos de regulação, avaliação e supervisão da educação superior no
    sistema federal de educação, e o Cadastro e-MEC de Instituições e Cursos
    Superiores e consolida disposições sobre indicadores de qualidade, banco de
    avaliadores (Basis) e o Exame Nacional de Desempenho de Estudantes (ENADE)
    e outras disposições.
    Disponível em: \url{http://portal.mec.gov.br/cne/arquivos/pdf/2007/port40_07.pdf}.
    Acesso em: 20 de janeiro de 2023.
    
    \item BRASIL. Comissão Nacional de Avaliação da Educação Superior.
    Resolução 1, de 17 de junho de 2010. Normatiza o Núcleo Docente
    Estruturante e dá outras providências. 
    Disponível em: \url{http://portal.mec.gov.br/index.php?option=com_docman&task=doc_download&gid=6885&Itemid}.
    Acesso em: 20 de janeiro de 2023.
    
    %\item BRASIL. Presidência da República. Casa Civil. Subchefia para Assuntos
    %Jurídicos. Decreto 5.622. Regulamenta o art. 80 da Lei 9.394, de 20 de
    %dezembro de 1996, que estabelece as diretrizes e bases da educação
    %nacional. 
    %Disponível em: \url{http://www.planalto.gov.br/ccivil_03/_ato2004-2006/2005/Decreto/D5622compilado.htm}.
    %%% REVOGADO
    \item BRASIL. Presidência da República. Casa Civil. Subchefia para Assuntos
    Jurídicos. Decreto 9.057. Regulamenta o art. 80 da Lei 9.394, de 20 de
    dezembro de 1996, que estabelece as diretrizes e bases da educação
    nacional. 
    Disponível em: \url{https://www.planalto.gov.br/ccivil_03/_Ato2015-2018/2017/Decreto/D9057.htm}.
    Acesso em: 20 de janeiro de 2023.
    
    \item BRASIL. Ministério da Educação. Portaria nº 4.059, de 10 de dezembro
    de 2004. Regulamentação de disciplinas na modalidade semipresencial.
    Disponível em: \url{http://portal.mec.gov.br/sesu/arquivos/pdf/nova/acs_portaria4059.pdf}.
    Acesso em: 20 de janeiro de 2023.
    
    \item BRASIL. Ministério da Educação. Conselho Nacional de Educação. Câmara
    de Educação Superior Parecer CNE/CES nº 136/2012, aprovado em 8 de março de
    2012 - Diretrizes Curriculares Nacionais para os cursos de graduação em
    Computação. 
    Publicado no DOU no 134, de 12 de julho de 2012.
    Disponível em: \url{http://portal.mec.gov.br/component/content/article?id=12991}.
    Acesso em: 20 de janeiro de 2023.
    
    \item BRASIL. Ministério da Educação. Conselho Nacional de Educação.
    Conselho Pleno. Parecer CNE/CP nº 003, de 10 de março de 2004.
    Disponível em: \url{http://portal.mec.gov.br/cne/arquivos/pdf/003.pdf}. 
    Acesso em: 20 de janeiro de 2023.
    
    \item BRASIL. Ministério da Educação. Conselho Nacional de Educação.
    Conselho Pleno. Resolução nº 1, de 30 de maio de 2012. Estabelece
    Diretrizes Nacionais para a Educação em Direitos Humanos. 
    Disponível em: \url{http://portal.mec.gov.br/index.php?option=com_docman&view=download&alias=10889-
    rcp001-12&category_slug=maio-2012-pdf&Itemid=30192}. 
    Acesso em: 20 de janeiro de 2023.
    
    \item BRASIL. Presidência da República. Casa Civil. Subchefia para Assuntos
    Jurídicos. Lei 12.764, de 27 de dezembro de 2012. Institui a Política
    Nacional de Proteção dos Direitos da Pessoa com Transtorno do Espectro
    Autista; e altera o § 3º do art. 98 da Lei 8.112, de 11 de dezembro de
    1990. 
    Disponível em: \url{http://www.planalto.gov.br/ccivil_03/_ato2011-2014/2012/lei/l12764.htm}.
    Acesso em: 20 de janeiro de 2023.
    
    \item FUNDAÇÃO UNIVERSIDADE FEDERAL DO ABC. Projeto Pedagógico Institucional. Santo
    André, 2017. Disponível em:
    \url{https://www.ufabc.edu.br/images/imagens_a_ufabc/projeto-pedagogico-institucional.pdf}.
    Acesso em: 20 de janeiro de 2023.
    
    \item FUNDAÇÃO UNIVERSIDADE FEDERAL DO ABC. Plano de Desenvolvimento
    Institucional. Santo André, 2013. Disponível em:
    \url{https://www.ufabc.edu.br/a-ufabc/documentos/plano-de-desenvolvimento-institucional-pdi}
    Acesso em: 20 de janeiro de 2023.
    
    \item BRASIL. Ministério da Educação. Conselho Nacional de Educação. Câmara
    de Educação Superior. Parecer CNE/CES nº 136/2012, aprovado em 8 de março
    de 2012 - Diretrizes Curriculares Nacionais para os cursos de graduação em
    Computação. 
    Disponível em: \url{http://portal.mec.gov.br/index.php?option=com_docman&view=download&alias=11205-pces136-11-pdf&category_slug=julho-2012-pdf&Itemid=30192}.
    Acesso em: 20 de janeiro de 2023.
    
    \item BRASIL. Ministério da Educação. Conselho Nacional de Educação. Câmara
    de Educação Superior. Resolução CNE/CES 5, de 16 de novembro de 2016 -
    Institui as Diretrizes Curriculares Nacionais para os cursos de graduação
    na área da Computação, abrangendo os cursos de bacharelado em Ciência da
    Computação, em Sistemas de Informação, em Engenharia de Computação, em
    Engenharia de Software e de licenciatura em Computação, e dá outras
    providências. Disponível em:
    \url{http://portal.mec.gov.br/index.php?option=com_docman&view=download&alias=52101-rces005-16-pdf&category_slug=novembro-2016-pdf&Itemid=30192}.
    Acesso em: 20 de janeiro de 2023.
    
\end{itemize}


\subsection{Regime de ensino}

\subsubsection{Estrutura curricular}

A estrutura do BCC é composta por três grupos de componentes curriculares, que
totalizam \red{3.246} horas: ($i$) disciplinas, ($ii$) atividades de extensão e
cultura e ($iii$) atividades complementares.
As disciplinas, por sua vez, estão divididas em três categorias: obrigatórias
(que devem necessariamente ser cursadas com aprovação para a integralização do
curso), de opção limitada (presentes em um grupo selecionado de disciplinas,
que permitem ao aluno aprofundar seus conhecimentos em determinadas áreas do
conhecimento, fazendo relações interdisciplinares com os conhecimentos
ofertados pelas disciplinas obrigatórias) e livres (quaisquer disciplinas
oferecidas pela UFABC, ou por outra IES reconhecida pelo MEC, de curso de
graduação ou de pós-graduação, necessárias para completar o número total de
créditos exigidos para a integralização do curso).
As disciplinas de opção limitada e livres, portanto, podem ser selecionadas
pelos estudantes, oferecendo autonomia para projetarem esta carga horária de
acordo com seus interesses e aptidões.
As disciplinas obrigatórias são apresentadas na
Seção~\ref{sec:matriz_curricular} e as disciplinas de opção limitada devem ser
selecionadas dentre aquelas constantes no Catálogo de Disciplinas de Opção
Limitada do BCC.

As componentes curriculares são contabilizadas na forma de créditos, sendo que
cada crédito equivale a 12 horas de aula e a hora-aula é de 60 minutos.
A distribuição da quantidade de créditos e da carga horária a serem cumpridos
em cada uma das categorias de disciplinas para a obtenção do grau de Bacharel
em Ciência da Computação é dada na Tabela~\ref{tab:carga_horaria}.
Não contabiliza-se, nessa tabela, o número de horas de estudo individual
extraclasse necessários para o bom acompanhamento das atividades.

\begin{table}[h!]
    \centering
    \caption{Distribuição de créditos e carga horária a serem cumpridos no BCC,
    descontando-se as horas de estudo individual extraclasse.}
    \label{tab:carga_horaria}
    \begin{tabular}{|l|c|c|}
        \hline
        \textbf{Categoria}                        & \textbf{Créditos} & \textbf{Carga horária (horas)} \\
        \hline\hline
        Disciplinas obrigatórias do BC\&T         & 84                & 1.008 \\
        \hline
        Disciplinas obrigatórias do BCC           & 108               & 1.296 \\
        \hline
        Disciplinas de opção limitada do BCC      & 24                & 288 \\
        \hline
        Disciplinas livres                        & 16                 & 192 \\
        \hline
        \textit{Total de disciplinas}             & \textit{232}       & \textit{2.784}\\
        \hline\hline
        \multicolumn{2}{|l|}{Trabalho de Conclusão de Curso}       & 288 \\
        \multicolumn{2}{|l|}{Atividades complementares}            & 18 \\
        \hline
        \multicolumn{2}{|l|}{Atividades de extensão e cultura}      & 343 \\
        \hline
        \multicolumn{2}{|l|}{\textit{Carga horária total}}               & \textbf{3.433}\\
        \hline
    \end{tabular}
\end{table}




\subsubsection{Interdisciplinaridade}

O BC\&T é a base da matriz curricular do BCC, de maneira que a formação
proposta proporciona interdisciplinaridade e flexibilidade curricular.
As disciplinas obrigatórias do BC\&T organizam o conhecimento em eixos
(Energia, Processos de Transformação, Representação e Simulação, Informação e
Comunicação, Estrutura da Matéria e Humanidades), visando despertar o interesse
dos alunos para a investigação de cunho interdisciplinar.
Os cursos de graduação da UFABC estão estruturados em um sistema de créditos
que permite diferentes organizações curriculares, de acordo com os interesses e
aptidões dos alunos.
Através das disciplinas livres, os alunos poderão se aprofundar em quaisquer
áreas do conhecimento, partindo para especificidades curriculares de cursos de
formação profissional ou explorando a interdisciplinaridade e estabelecendo um
currículo individual de formação.

É importante destacar que a interdisciplinaridade do presente projeto
pedagógico e a possibilidade de escolher disciplinas livres permitem que o
discente formado no BCC da UFABC esteja alinhado com as seguintes diretrizes
legais:
\begin{itemize}
    \item Decreto 5.626, de 22 de dezembro de 2005: a disciplina de LIBRAS,
    cuja ementa faz parte do rol de disciplinas dos cursos de licenciatura da
    UFABC, pode ser cursada pelos alunos do BCC.
    
    \item Lei nº 11.64, sobre a obrigatoriedade da temática ``História e
    Cultura Afro-Brasileira e Indígena'', e Resolução 01/2004, de 17 de junho
    de 2004: o aluno do BCC pode escolher cursar disciplinas livres que fazem
    parte do rol de disciplinas da UFABC e que envolvem a temática da História
    e Cultura Afro-Brasileira e Indígenas.
    
    \item Política Nacional de Educação Ambiental (Lei nº 9795/1999 e decreto
    4.281, de 25 de junho de 2002): muitas disciplinas livres oferecidas no rol
    de disciplinas de engenharia ambiental podem ser cursadas pelos alunos do
    BCC, permitindo assim a integração desse projeto pedagógico com a educação
    ambiental.
\end{itemize}



\subsection{Estratégias pedagógicas}

A UFABC foi concebida definindo a interdisciplinaridade como uma referência
pedagógica.
Sendo assim, o desenvolvimento do perfil do egresso é trabalhado com uma
formação interdisciplinar com alto grau de liberdade para incorporar
componentes curriculares à sua formação.
Além de cobrir os assuntos pertinentes à formação definida pelas DCNs de
Computação, esse modelo possibilita que o aluno desenvolva competências em
outras áreas de seu próprio interesse.

Seguindo a recomendação da matriz curricular, os primeiros quadrimestres
letivos de curso são preenchidos por disciplinas do BC\&T, onde o(a) aluno(a)
tem contato com várias áreas da Ciência, fortalecendo sua base científica e
humanística, além de experimentar os primeiros contatos com disciplinas da área
de Computação (Bases Computacionais da Ciência, Natureza da Informação,
Comunicação e Redes, Processamento da Informação).
Aos poucos, o(a) aluno(a) vai encontrar janelas de horários para incluir
disciplinas específicas de Computação enquanto finaliza sua formação no BC\&T.
O projeto pedagógico ainda prevê aproximadamente \red`{15\%} de carga horária exclusivamente 
em disciplinas livres e de opção limitada, em que o(a) aluno(a) poderá escolher os componentes
curriculares que completarão a sua formação.

Na UFABC as disciplinas não possuem pré-requisitos entre si.
Mesmo assim, a estrutura da matriz curricular sugere uma composição que
favorece o desenvolvimento contínuo das competências e habilidades do egresso
durante o desenvolvimento do curso, concentrando disciplinas que abordam temas
avançados e específicos no final do curso e disciplinas fundamentais em seu
início.

O uso de Tecnologias de Informação e Comunicação (TICs) é uma realidade próxima
dos estudantes na UFABC.
Muitas disciplinas utilizam ambientes virtuais de aprendizagem (AVAs) para
gestão de conteúdo em disciplinas presenciais e semipresenciais.
Durante a pandemia de COVID-19, algumas disciplinas foram também ofertadas na
modalidade online com sucesso.
Todos os cursos possuem páginas específicas em que seus conteúdos e documentos
ficam acessíveis à comunidade (projeto pedagógico, informações gerais,
documentos, links para outras páginas de recursos, etc.).
Uma importante parcela da carga horária total é trabalhada em aulas práticas,
ofertadas em laboratórios de informática com computadores ou laboratórios de
{\it hardware} com dispositivos eletrônicos.

Em termos de acessibilidade, a UFABC tem se preocupa com a garantia de acesso
às pessoas com deficiência e/ou com mobilidade reduzida. 
Seguindo as determinações do Decreto nº 5.296/2004 47 e da Lei 10.098/2000 48,
os dois campi da UFABC possuem acessibilidade arquitetônica, garantindo o uso
autônomo dos espaços por pessoas com deficiência e/ou com mobilidade reduzida.
Através do Núcleo de Acessibilidade da Pró-Reitoria de Assuntos Comunitários e
Políticas Afirmativas (PROAP), a UFABC tem procurado a excelência no quesito
inclusão.
Nesse sentido, dentre as disciplinas oferecidas pela UFABC, destacamos o
oferecimento da disciplina Libras.

Políticas de educação ambiental e de educação em direitos humanos são tratadas
por algumas disciplinas ofertadas pela UFABC.
Dentre as relacionadas à educação ambiental, citamos: Educação Ambiental,
Economia e Meio Ambiente; Economia, Sociedade e Meio Ambiente; e Energia, Meio
Ambiente e Sociedade.
Dentre as relacionadas à educação em direitos humanos, citamos: Regime
Internacional dos Direitos Humanos e a Atuação Brasileira; Movimentos
Sindicais, Sociais e Culturais; Diversidade Cultural, Conhecimento Local e
Políticas Públicas; Identidade e Cultura; Cidadania, Direitos e Desigualdades;
Cultura, Identidade e Política na América Latina; e Trajetórias Internacionais
do Continente Africano.



\subsection{Matriz Curricular Recomendada}
\label{sec:matriz_curricular}

Os componentes curriculares na UFABC são oferecidos em ciclos quadrimestrais.
A carga horária das disciplinas são distribuídas entre créditos teóricos (T),
práticos (P), de caráter extensionista ou cultural (E) e de dedicação a estudos
individuais extraclasses (I).
Considera-se, dessa forma, a quantidade de créditos e de horas de trabalho de
cada disciplina apresentada por seu T-P-E-I.
Sugere-se que o aluno pondere o número de horas de estudo individual
extraclasse nos momentos de matrícula, para que considere sempre as horas
necessárias de dedicação às atividades de cada disciplina. 
Para o cômputo dos créditos totais, no entanto, são considerados apenas os
especificados em T e P e, no caso de E, serão contabilizados para compor a
carga horária extensionista e cultural.


Os estudantes da UFABC têm liberdade para organização da própria matriz
curricular, não havendo pré-requisitos entre as disciplinas.
No entanto, é importante observar os requisitos recomendados e o encadeamento
das disciplinas de modo a permitir o melhor aproveitamento dos conteúdos.
Com o objetivo de orientar e auxiliar o aluno a compreender as possibilidades desse
currículo, a Figura~\ref{fig:matriz_curricular} apresenta a Matriz Curricular
Sugerida para o discente que pretende integralizar ambos BC\&T e o BCC em cinco
anos e a Figura~\ref{fig:recomendacoes} apresenta graficamente as recomendações
entre as disciplinas obrigatórias do BCC e do BC\&T.
A descrição completa das ementas das disciplinas obrigatórias encontra-se na
Seção~\ref{sec:disciplinas_obrigatorias}.

\begin{figure}
    \centering
    \resizebox{\textwidth}{!}{
\documentclass{standalone}
\usepackage{amsmath, amsthm, amsfonts, amssymb}
\usepackage{tikz}
\usetikzlibrary{shapes,snakes,positioning,calc}
\usepackage{hyperref}

\newcommand{\disciplina}[6][nyellow]{
\node [draw=black, fill=#1, very thick, rectangle, rounded corners, inner
sep=2pt, inner ysep=2pt, text width=#5pt, minimum height=40pt,
text centered, anchor=west] (#4) at (#2,#3) {
   \begin{minipage}{#5pt}
   \linespread{1.0}\selectfont
       \centering
       \footnotesize{#6}
   \end{minipage}
};
}


\usepackage{xcolor}
\definecolor{nred}{rgb}{0.88, 0.28, 0.33}
\definecolor{nblue}{rgb}{0.34, 0.74, 0.96}
\definecolor{nyellow}{rgb}{0.8, 0.86, 0.38}
\definecolor{npurple}{rgb}{0.25, 0.34, 0.93}
\definecolor{ngreen}{rgb}{0.45, 0.95, 0.66}


\begin{document}

\begin{tikzpicture}

    \def\sz{22pt}
    \def\doisc{38} % 2*\sz - 6
    \def\tresc{60} % 3*\sz - 6
    \def\quatc{82} % 4*\sz - 6
    \def\cincc{104} % 5*\sz - 6
    \def\seisc{126} % 6*\sz - 6
    
    %Q1
    \node [draw, rotate=90, black,rectangle, minimum width=120pt, minimum
    height=10pt, rounded corners] at (-25pt,-40pt) {\footnotesize{1º ANO}};
    
    \node [draw, rotate=90, black,rectangle, minimum width=40pt, minimum
    height=10pt, rounded corners] at (-10pt,0pt) {\footnotesize{1º quad}};
    
    \disciplina [gray!60]{0*\sz}{0pt}{a}{\tresc}{Humanidades 1 \\(3-0-0-4)}
    
    \disciplina [gray!60]{3*\sz}{0pt}{b}{\tresc}{Base Experimental das Ciências Naturais\\ (0-3-0-2)} 
    
    \disciplina [gray!60]{6*\sz}{0pt}{c}{\tresc}{Estrutura da Matéria \\(3-0-0-4)}
    
    \disciplina [gray!60]{9*\sz}{0pt}{d}{\quatc}{Bases Matemáticas\\(4-0-0-5)}
    
    \disciplina [gray!60]{13*\sz}{0pt}{e}{\tresc}{Evolução e Diversificação da Vida na Terra \\(3-0-0-4)}
    
    \disciplina [gray!60]{16*\sz}{0pt}{f}{\doisc}{\scriptsize Bases Computacionais da Ciência \\(0-2-0-2)}
    
    
    %Q2
    \node [draw, rotate=90, black,rectangle, minimum width=40pt, minimum
    height=10pt, rounded corners] at (-10pt,-40pt) {\footnotesize{2º quad}};
    
    \disciplina [gray!60]{0*\sz}{-40pt}{a}{\tresc}{Natureza da Informação\\(3-0-0-4)}
    
    \disciplina [gray!60]{3*\sz}{-40pt}{a}{\cincc}{Fenômenos Mecânicos\\(4-1-0-6)}
    
    \disciplina [gray!60]{8*\sz}{-40pt}{a}{\quatc}{Funções de uma Variável\\(4-0-0-6)}
    
    \disciplina [gray!60]{12*\sz}{-40pt}{a}{\tresc}{Geometria Analítica\\(3-0-0-6)}
    
    \disciplina [gray!60]{15*\sz}{-40pt}{a}{\tresc}{Biodiversidade: Int. Organismos e Ambiente\\(3-0-0-4)}
    
    
    %Q3
    \node [draw, rotate=90, black,rectangle, minimum width=40pt, minimum
    height=10pt, rounded corners] at (-10pt,-80pt) {\footnotesize{3º quad}};
    
    \disciplina [gray!60]{0*\sz}{-80pt}{a}{\quatc}{Processamento da Informação\\(0-4-0-4)}
    
    \disciplina [gray!60]{4*\sz}{-80pt}{a}{\cincc}{Fenômenos Térmicos\\(4-1-0-6)}
    
    \disciplina [gray!60]{9*\sz}{-80pt}{a}{\quatc}{Funções de Várias Variáveis\\(4-0-0-6)}
    
    \disciplina [gray!60]{13*\sz}{-80pt}{a}{\cincc}{Transformações Químicas\\(3-2-0-6)}
    
    
    %Q4
    \node [draw, rotate=90, black,rectangle, minimum width=120pt, minimum
    height=10pt, rounded corners] at (-25pt,-160pt) {\footnotesize{2º ANO}};
    
    \node [draw, rotate=90, black,rectangle, minimum width=40pt, minimum
    height=10pt, rounded corners] at (-10pt,-120pt) {\footnotesize{4º quad}};
    
    \disciplina [gray!60]{0*\sz}{-120pt}{a}{\tresc}{Comunicação e Redes\\(3-0-0-4)}
    
    \disciplina [gray!60]{3*\sz}{-120pt}{a}{\cincc}{Fenômenos Eletromagnéticos\\(4-1-0-6)}
    
    \disciplina [gray!60]{8*\sz}{-120pt}{a}{\quatc}{Introdução às Equações Diferencias Ordinárias\\(4-0-0-4)}
    
    \disciplina [gray!60]{12*\sz}{-120pt}{a}{\tresc}{Introdução à Probabilidade e Estatística\\(3-0-0-4)}
    
    \disciplina [nblue]{15*\sz}{-120pt}{a}{\quatc}{\hyperref[disc:mdI]{Matemática Discreta I}\\(4-0-0-4)}

    
    %Q5
    \node [draw, rotate=90, black,rectangle, minimum width=40pt, minimum
    height=10pt, rounded corners] at (-10pt,-160pt) {\footnotesize{5º quad}};
    
    \disciplina [gray!60]{0*\sz}{-160pt}{a}{\tresc}{Humanidades 2\\(3-0-0-4)}
    
    \disciplina [gray!60]{3*\sz}{-160pt}{a}{\cincc}{Bioquímica: Estrutura, Propriedade e Funções de Biomoléculas\\(3-2-0-6)}
    
    \disciplina [gray!60]{8*\sz}{-160pt}{a}{\quatc}{Física Quântica\\(4-0-0-4)}
    
    \disciplina [nblue]{12*\sz}{-160pt}{a}{\quatc}{\hyperref[disc:pe]{Programação Estruturada}\\(2-2-0-6)}
    

    %Q6
    \node [draw, rotate=90, black,rectangle, minimum width=40pt, minimum
    height=10pt, rounded corners] at (-10pt,-200pt) {\footnotesize{6º quad}};
    
    \disciplina [gray!60]{0*\sz}{-200pt}{a}{\tresc}{Humanidades 3\\(3-0-0-4)}
    
    \disciplina [nblue]{3*\sz}{-200pt}{a}{\quatc}{\hyperref[disc:circ_dig]{Circuitos Digitais}\\(3-1-0-4)}
    
    \disciplina [nblue]{7*\sz}{-200pt}{a}{\seisc}{\hyperref[disc:alge_lin]{Álgebra Linear}\\(6-0-0-5)}
    
    \disciplina [nblue]{13*\sz}{-200pt}{a}{\quatc}{\hyperref[disc:aedI]{Algoritmos e Estruturas de Dados I}\\(2-2-0-6)}
    

    %Q7
    \node [draw, rotate=90, black,rectangle, minimum width=120pt, minimum
    height=10pt, rounded corners] at (-25pt,-280pt) {\footnotesize{3º ANO}};
    
    \node [draw, rotate=90, black,rectangle, minimum width=40pt, minimum
    height=10pt, rounded corners] at (-10pt,-240pt) {\footnotesize{7º quad}};
    
    \disciplina [nblue]{0*\sz}{-240pt}{a}{\quatc}{\hyperref[disc:poo]{Programação Orientada a Objetos}\\(2-2-0-4)}
    
    \disciplina [nblue]{4*\sz}{-240pt}{a}{\quatc}{\hyperref[disc:sist_dig]{Sistemas Digitais}\\(2-2-0-4)}
    
    \disciplina [nblue]{8*\sz}{-240pt}{a}{\quatc}{\hyperref[disc:mdII]{Matemática Discreta II}\\(4-0-0-4)}
    
    \disciplina [nblue]{12*\sz}{-240pt}{a}{\quatc}{\hyperref[disc:aedII]{Algoritmos e Estruturas de Dados II}\\(4-0-0-6)}
    
    
    %Q8
    \node [draw, rotate=90, black,rectangle, minimum width=40pt, minimum
    height=10pt, rounded corners] at (-10pt,-280pt) {\footnotesize{8º quad}};
    
    \disciplina [nyellow]{0*\sz}{-280pt}{a}{\quatc}{Opção Limitada\\(4 créditos)}
    
    \disciplina [nblue]{4*\sz}{-280pt}{a}{\quatc}{\hyperref[disc:arq]{Arquitetura de Computadores}\\(4-0-0-4)}
    
    \disciplina [nblue]{8*\sz}{-280pt}{a}{\quatc}{\hyperref[disc:ia]{Inteligência Artificial}\\(4-0-0-4)}
    
    \disciplina [nblue]{12*\sz}{-280pt}{a}{\quatc}{\hyperref[disc:ag]{Algoritmos em Grafos}\\(4-0-0-4)}
    
    
    %Q9
    \node [draw, rotate=90, black,rectangle, minimum width=40pt, minimum
    height=10pt, rounded corners] at (-10pt,-320pt) {\footnotesize{9º quad}};
    
    \disciplina [nyellow]{0*\sz}{-320pt}{a}{\quatc}{Opção Limitada\\(4 créditos)}
    
    \disciplina [nblue]{4*\sz}{-320pt}{a}{\quatc}{\hyperref[disc:redes]{Redes de Computadores}\\(3-1-0-4)}
    
    \disciplina [nblue]{8*\sz}{-320pt}{a}{\quatc}{\hyperref[disc:es]{Engenharia de Software}\\(4-0-0-4)}
    
    \disciplina [nblue]{12*\sz}{-320pt}{a}{\quatc}{\hyperref[disc:aaI]{Análise de Algoritmos I}\\(4-0-0-4)}
    

    %Q10
    \node [draw, rotate=90, black,rectangle, minimum width=120pt, minimum
    height=10pt, rounded corners] at (-25pt,-400pt) {\footnotesize{4º ANO}};
    
    \node [draw, rotate=90, black,rectangle, minimum width=40pt, minimum
    height=10pt, rounded corners] at (-10pt,-360pt) {\footnotesize{10º quad}};
    
    \disciplina [nyellow]{0*\sz}{-360pt}{a}{\quatc}{Opção Limitada\\(4 créditos)}
    
    \disciplina [nblue]{4*\sz}{-360pt}{a}{\doisc}{\hyperref[disc:ces]{\scriptsize Computadores, Ética e Sociedade}\\(2-0-0-4)}
    
    \disciplina [nblue]{6*\sz}{-360pt}{a}{\quatc}{\hyperref[disc:pf]{Programação Funcional}\\(4-0-0-4)}
    
    \disciplina [nblue]{10*\sz}{-360pt}{a}{\quatc}{\hyperref[disc:aaII]{Análise de Algoritmos II}\\(4-0-0-4))}
    
    
    %Q11
    \node [draw, rotate=90, black,rectangle, minimum width=40pt, minimum
    height=10pt, rounded corners] at (-10pt,-400pt) {\footnotesize{11º quad}};
    
    \disciplina [nyellow]{0*\sz}{-400pt}{a}{\quatc}{Opção Limitada\\(4 créditos)}
    
    \disciplina [nblue]{4*\sz}{-400pt}{a}{\quatc}{\hyperref[disc:mbd]{Modelagem de Banco de Dados}\\(4-0-0-4)}
    
    \disciplina [nblue]{8*\sz}{-400pt}{a}{\quatc}{\hyperref[disc:cg]{Computação Gráfica}\\(0-4-0-4)}
    
    \disciplina [nblue]{12*\sz}{-400pt}{a}{\quatc}{\hyperref[disc:ol]{Otimização Linear}\\(4-0-0-4)}
    
    %Q12
    \node [draw, rotate=90, black,rectangle, minimum width=40pt, minimum
    height=10pt, rounded corners] at (-10pt,-440pt) {\footnotesize{12º quad}};
    
    \disciplina [nyellow]{0*\sz}{-440pt}{a}{\quatc}{Opção Limitada\\(4 créditos)}
    
    \disciplina [nblue]{4*\sz}{-440pt}{a}{\doisc}{\scriptsize Metod.\ e Escr.\ Cient.\ para Comp.\\(2-0-0-2)}
    
    \disciplina [nblue]{6*\sz}{-440pt}{a}{\quatc}{\hyperref[disc:so]{Sistemas Operacionais}\\(3-1-0-4)}
    
    \disciplina [nblue]{10*\sz}{-440pt}{a}{\quatc}{\hyperref[disc:lfa]{Linguagens Formais e Autômatos}\\(4-0-0-4)}
    

    %Q13
    \node [draw, rotate=90, black,rectangle, minimum width=120pt, minimum
    height=10pt, rounded corners] at (-25pt,-520pt) {\footnotesize{5º ANO}};
    
    \node [draw, rotate=90, black,rectangle, minimum width=40pt, minimum
    height=10pt, rounded corners] at (-10pt,-480pt) {\footnotesize{13º quad}};
    
    \disciplina [nyellow]{0*\sz}{-480pt}{a}{\quatc}{Opção Limitada\\(4 créditos)}
    
    \disciplina [nblue]{4*\sz}{-480pt}{a}{\quatc}{Trabalho de Conclusão de Curso I\\(4-0-0-4)}
    
    \disciplina [nblue]{8*\sz}{-480pt}{a}{\quatc}{\hyperref[disc:compi]{Compiladores e Interpretadores}\\(4-0-0-4)}
    
    
    %Q14
    \node [draw, rotate=90, black,rectangle, minimum width=40pt, minimum
    height=10pt, rounded corners] at (-10pt,-520pt) {\footnotesize{14º quad}};
    
    \disciplina [nred]{0*\sz}{-520pt}{a}{\quatc}{Livre\\(4 créditos)}
    
    \disciplina [nblue]{4*\sz}{-520pt}{a}{\quatc}{Trabalho de Conclusão de Curso II\\(4-0-0-6)}
    
    \disciplina [nblue]{8*\sz}{-520pt}{a}{\quatc}{\hyperref[disc:sist_dist]{Sistemas Distribuídos}\\(3-1-0-4)}
    

    %Q15
    \node [draw, rotate=90, black,rectangle, minimum width=40pt, minimum
    height=10pt, rounded corners] at (-10pt,-560pt) {\footnotesize{15º quad}};
    
    \disciplina [nred]{0*\sz}{-560pt}{a}{\quatc}{Livre\\(4 créditos)}

    \disciplina [nblue]{4*\sz}{-560pt}{a}{\quatc}{Trabalho de Conclusão de Curso III\\(4-0-0-6)}
    
    \disciplina [nblue]{8*\sz}{-560pt}{a}{\quatc}{\hyperref[disc:seg]{Segurança de Dados}\\(3-1-0-4)}
    
    
    \draw [red,thick,dashed] (20*\sz,20pt) -- (20*\sz,-580pt); 
    \node [text=red] at (420pt,35pt) {20 créditos};


    % Legenda
    \node [text=black] at (-20pt, -620pt) {\textbf{Legenda:}};
    \disciplina [gray!60]{15pt}{-620pt}{bis0505}{63}{Obrigatória BC\&T\\(T-P-E-I)}
    
    \disciplina [nblue]{88pt}{-620pt}{bis0505}{63}{Obrigatória BCC\\(T-P-E-I)}
    
    \disciplina [nyellow]{161pt}{-620pt}{bis0505}{63}{Opção Limitada BCC\\(T-P-E-I)}
    
    \disciplina [nred]{234pt}{-620pt}{bis0505}{63}{Livre BCC\\(T-P-E-I)}
    
    
\end{tikzpicture}


\end{document}
}
    \caption{Matriz curricular recomendada para integralização do BC\&T e do
    BCC em 5 anos. A largura dos retângulos é diretamente proporcional à
    quantidade de créditos em sala de aula.}
    \label{fig:matriz_curricular}
\end{figure}

\begin{figure}
    \centering
    \resizebox{\textwidth}{!}{
\documentclass{standalone}
\usepackage{amsmath, amsthm, amsfonts, amssymb}
\usepackage{tikz}
\usetikzlibrary{shapes,snakes,positioning,calc}

\newcommand{\disciplina}[6][nyellow]{
\node [draw=black, fill=#1, very thick, rectangle, rounded corners, inner
sep=2pt, inner ysep=2pt, text width=#5pt, minimum height=40pt,
text centered, anchor=west] (#4) at (#2,#3) {
   \begin{minipage}{#5pt}
   \linespread{1.0}\selectfont
       \centering
       \scriptsize{#6}
   \end{minipage}
};
}
\newcommand{\disciplinab}[6][yellow!20]{
\node [draw=black, fill=#1, very thick, rectangle, rounded corners, inner
sep=2pt, inner ysep=2pt, text width=#5pt, minimum height=40pt,
text centered, anchor=west] (#4) at (#2,#3) {
   \begin{minipage}{#5pt}
   \linespread{1.0}\selectfont
       \centering
       \small{#6}
   \end{minipage}
};
}


\usepackage{xcolor}
\definecolor{nred}{rgb}{0.88, 0.28, 0.33}
\definecolor{nblue}{rgb}{0.34, 0.74, 0.96}
\definecolor{nyellow}{rgb}{0.8, 0.86, 0.38}
\definecolor{npurple}{rgb}{0.25, 0.34, 0.93}
\definecolor{ngreen}{rgb}{0.45, 0.95, 0.66}


\begin{document}

\begin{tikzpicture}

    \def\sz{21pt}
    \def\doisc{36} % 2*\sz - 6
    \def\tresc{57} % 3*\sz - 6
    \def\quatc{78} % 4*\sz - 6
    \def\cincc{99} % 5*\sz - 6
    \def\seisc{120} % 6*\sz - 6
    \def\h{40pt}


    \disciplinab [nblue]{-1*\sz}{0*\h}{PE}{\tresc}{Programação Estruturada}
    \disciplinab [nblue]{-1*\sz}{-2*\h}{AEDI}{\tresc}{Algoritmos e Estruturas de Dados I}
    \disciplinab [nblue]{-1*\sz}{-4*\h}{AEDII}{\tresc}{Algoritmos e Estruturas de Dados II}
    \disciplinab [nblue]{-1*\sz}{-6*\h}{AG}{\tresc}{Algoritmos em Grafos}
    \disciplinab [nblue]{-1*\sz}{-8*\h}{AAI}{\tresc}{Análise de Algoritmos I}
    \disciplinab [nblue]{-1*\sz}{-10*\h}{AAII}{\tresc}{Análise de Algoritmos II}
    \disciplinab [nblue]{-1*\sz}{-12*\h}{LFA}{\tresc}{Linguagens Formais e Autômatos}
    \disciplinab [nblue]{-1*\sz}{-14*\h}{COMP}{\tresc}{Compiladores e Interpretadores}

    \disciplinab [nblue]{16*\sz}{-1*\h}{CG}{\tresc}{Computação Gráfica}
    \disciplinab [nblue]{5*\sz}{-1*\h}{IHC}{\tresc}{Interação Humano-Computador}
%    \disciplinab [nblue]{5*\sz}{-1*\h}{CG}{\tresc}{Computação Gráfica}
    \disciplinab [nblue]{5*\sz}{-3*\h}{PF}{\tresc}{Programação Funcional}
    \disciplinab [nblue]{5*\sz}{-5*\h}{IA}{\tresc}{Inteligência Artificial}
    \disciplinab [nblue]{5*\sz}{-8*\h}{ARQ}{\tresc}{Arquitetura de Computadores}
    \disciplinab [nblue]{5*\sz}{-10*\h}{SO}{\tresc}{Sistemas Operacionais}
    \disciplinab [nblue]{5*\sz}{-12*\h}{MDII}{\tresc}{Matemática Discreta II}

    \disciplinab [nblue]{10*\sz}{-4*\h}{CD}{\tresc}{Circuitos Digitais}
    \disciplinab [nblue]{10*\sz}{-6*\h}{SD}{\tresc}{Sistemas Digitais}
    \disciplinab [nblue]{10*\sz}{-10*\h}{MD}{\tresc}{Matemática Discreta I}

    \disciplinab [nblue]{15*\sz}{-14*\h}{AL}{\tresc}{Álgebra Linear}

    \disciplinab [nblue]{20*\sz}{0*\h}{POO}{\tresc}{Programação Orientada a Objetos}
    \disciplinab [nblue]{20*\sz}{-2*\h}{RD}{\tresc}{Redes de Computadores}
    \disciplinab [nblue]{20*\sz}{-4*\h}{SDT}{\tresc}{Sistemas Distribuídos}

    \disciplinab [nblue]{25*\sz}{0*\h}{ES}{\tresc}{Engenharia de Software}
    \disciplinab [nblue]{25*\sz}{-2*\h}{BD}{\tresc}{Modelagem de Banco de Dados}
    \disciplinab [nblue]{25*\sz}{-4*\h}{SEG}{\tresc}{Segurança de Dados}

    \disciplinab [nblue]{20*\sz}{-14*\h}{CES}{\tresc}{Computadores, Ética e Sociedade}
    \disciplinab [nblue]{25*\sz}{-14*\h}{MET}{\tresc}{Metod. e Escr. Cient. para Comp.}

    \disciplinab [gray!50]{-1*\sz}{2*\h}{PI}{\tresc}{Processamento da Informação}
    \disciplinab [gray!50]{5*\sz}{2*\h}{BCC}{\tresc}{Bases Computacionais da Ciência}
    \disciplinab [gray!50]{10*\sz}{2*\h}{NI}{\tresc}{Natureza da Informação}
    \disciplinab [gray!50]{15*\sz}{-4*\h}{FEMAG}{\tresc}{Fenômenos Eletromagnéticos}
    \disciplinab [gray!50]{15*\sz}{-8*\h}{FEMEC}{\tresc}{Fenômenos Mecânicos}
    \disciplinab [gray!50]{15*\sz}{-10*\h}{FUV}{\tresc}{Funções de Uma Variável}
    \disciplinab [gray!50]{20*\sz}{-10*\h}{GA}{\tresc}{Geometria Analítica}
    \disciplinab [gray!50]{15*\sz}{-12*\h}{BM}{\tresc}{Bases Matemáticas}
    \disciplinab [gray!50]{10*\sz}{-12*\h}{FVV}{\tresc}{Funções de Várias Variáveis}

    \draw[seta] (AEDI) -- (PE);
    \draw[seta] (AEDII) -- (AEDI);
    \draw[seta] (AG) -- (AEDII);
    \draw[seta] (AAI) -- (AG);
    \draw[seta] (AAII) -- (AAI);
    \draw[seta] (LFA) -- (AAII);
    \draw[seta] (AG) -- (MDII);
    \draw[seta] (MDII) -- (MD);
    %\draw[seta] (OL) -- (AAII);
    %\draw[seta] (OL) -- (AL);
    \draw[seta] (COMP) -- (LFA);
    \draw[seta] (CG) -- (AEDI);
    \draw[seta] (CG) -- (POO);
    \draw[seta] (POO) -- (PE);
    \draw[seta] (PF) -- (AEDI);
    \draw[seta] (IA) -- (AEDI);
    \draw[seta] (ARQ) -- (AEDI);
    \draw[seta] (ARQ) -- (SD);
    \draw[seta] (SD) -- (CD);
    \draw[seta] (SO) -- (ARQ);
    \draw[seta] (SO) -- (AEDII);
    \draw[seta] (ES) -- (POO);
    \draw[seta] (BD) -- (POO);
    \draw[seta] (RD) -- (POO);
    \draw[seta] (RD) -- (AEDI);
    \draw[seta] (SDT) -- (RD);
    \draw[seta] (SDT) -- (SO);
    \draw[seta] (SEG) -- (RD);

    \draw[seta,gray] (PE) -- (PI);
    \draw[seta,gray] (PI) -- (BCC);
    \draw[seta,gray] (NI) -- (BCC);
    \draw[seta,gray] (CD) -- (NI);
    \draw[seta,gray] (CD) -- (FEMAG);
    \draw[seta,gray] (FEMAG) -- (FEMEC);
    \draw[seta,gray] (FEMEC) -- (GA);
    \draw[seta,gray] (FEMEC) -- (FUV);
    \draw[seta,gray] (GA) -- (BM);
    \draw[seta,gray] (CG) -- (GA);
    \draw[seta,gray] (FUV) -- (BM);
    \draw[seta,gray] (AL) -- (GA);
    \draw[seta,gray] (MD) -- (FUV);
    %\draw[seta,gray] (OL) -- (FVV);
    \draw[seta,gray] (FVV) -- (FUV);
    \draw[seta,gray] (IHC) -- (PI);

\end{tikzpicture}


\end{document}
}
    \caption{Recomendações entre as disciplinas obrigatórias do BCC e do BC\&T.}
    \label{fig:recomendacoes}
\end{figure}

As regras de transição entre a matriz sugerida do PPC anterior e a matriz
sugerida do PPC atual são apresentadas na Seção~\ref{sec:regras_transicao}.

\subsection{Mapeamento de habilidades/competências e atividades pedagógicas}

A organização curricular foi desenhada para atender aos requisitos estruturais
e pedagógicos da UFABC, bem como às Diretrizes Curriculares Nacionais dos
cursos de graduação em Computação (Resolução CNE/CES nº 5, de 16 de novembro de
2016).
Na Tabela~\ref{tab:mapeamento_competencias}, indicamos os componentes
pedagógicos que contribuem para a formação e consolidação das habilidades e
competências dos egressos.
As atividades pedagógicas estão classificadas da seguinte forma:

\begin{itemize}
    \item \textcolor{nred}{Disciplinas obrigatórias do BC\&T}
    \item \textcolor{nblue}{Disciplinas obrigatórias do BCC}
    \item \textcolor{nyellow}{Disciplinas de opção limitada do BCC}
    \item \textcolor{npurple}{Outras ações}
\end{itemize}


\begin{longtable}{|p{.3\textwidth}p{.3\textwidth}p{.25\textwidth}|}
    \caption{Componentes pedagógicos e suas contribuições de acordo com as habilidades e competências dos egressos.}
    \label{tab:mapeamento_competencias}
    \endfirsthead
    \endhead

    \multicolumn{3}{p{0.95\textwidth}}{Identificar problemas que tenham solução algorítmica}\\
    \hline
    \textcolor{nred}{Bases Comput. da Ciência}  &
    \textcolor{nblue}{Algs. e Estruturas de Dados I} &
    \textcolor{nblue}{Otimização Linear} \\
    \textcolor{nred}{Processamento da Informação} &
    \textcolor{nblue}{Algs. e Estruturas de Dados II} &
    \textcolor{nblue}{Algs. em Grafos} \\
    \textcolor{nblue}{Progr. Estruturada} &
    \textcolor{nblue}{Análise de Algoritmos I} &
    \textcolor{nblue}{Inteligência Artificial} \\
    \textcolor{nblue}{Matemática Discreta} &
    \textcolor{nblue}{Análise de Algoritmos II} &
    \textcolor{nblue}{Progr. Funcional} \\
    \textcolor{nblue}{Matemática Discreta II} &
    \textcolor{nblue}{Progr. Orientada a Objetos} & \\

    \hline
    \multicolumn{3}{p{0.95\textwidth}}{}\\
    
    \multicolumn{3}{p{0.95\textwidth}}{Conhecer os limites da computação}\\
    \hline
    \textcolor{nblue}{Análise de Algoritmos I} &
    \textcolor{nblue}{Análise de Algoritmos II} &
    \textcolor{nblue}{Ling. Formais e Autômatos} \\

    \hline
    \multicolumn{3}{p{0.95\textwidth}}{}\\

    \multicolumn{3}{p{0.95\textwidth}}{Resolver problemas usando ambientes de programação}\\
    \hline
    \textcolor{nred}{Processamento da Informação} &
    \textcolor{nblue}{Algs. em Grafos} &
    \textcolor{nblue}{Engenharia de Software} \\
    \textcolor{nblue}{Progr. Estruturada} &
    \textcolor{nblue}{Inteligência Artificial} &
    \textcolor{nblue}{Otimização Linear}\\
    \textcolor{nblue}{Algs. e Estruturas de Dados I} &
    \textcolor{nblue}{Progr. Orientada a Objetos} &
    \textcolor{nblue}{Progr. Funcional}\\
    \textcolor{nblue}{Algs. e Estruturas de Dados II} &
    \textcolor{nblue}{\small Model. de Banco de Dados} &
    \textcolor{nblue}{Sistemas Digitais}\\
    \textcolor{nblue}{Compiladores e Interpretadores} & & \\
    \hline
    
    \multicolumn{3}{p{0.95\textwidth}}{}\\

    \multicolumn{3}{p{0.95\textwidth}}{Tomar decisões e inovar, com base no
    conhecimento do funcionamento e das características técnicas de hardware
    e da infraestrutura de software dos sistemas de computação consciente dos
    aspectos éticos, legais e dos impactos ambientais decorrentes}\\
    \hline
    \textcolor{nred}{Ciência, Tecnologia e Sociedade} &
    \textcolor{nblue}{Redes de Computadores} &
    \textcolor{nblue}{Segurança de Dados}\\
    \textcolor{nred}{Comunicação e Redes} &
    \textcolor{nblue}{Sistemas Operacionais} &
    \textcolor{nblue}{\small Model. de Banco de Dados}\\
    \textcolor{nblue}{Arquitetura de Computadores} &
    \textcolor{nblue}{Sistemas Distribuídos} &
    \textcolor{nblue}{Engenharia de Software} \\
    \textcolor{nblue}{Sistemas Digitais} &
    \textcolor{nblue}{Comput., Ética e Sociedade} & \\
    \hline
    
    \multicolumn{3}{p{0.95\textwidth}}{}\\
    \multicolumn{3}{p{0.95\textwidth}}{Compreender e explicar as dimensões quantitativas de um
    problema}\\
    \hline
    \textcolor{nred}{Natureza da Informação} &
    \textcolor{nred}{Geometria Analítica} &
    \textcolor{nblue}{Análise de Algoritmos I} \\
    \textcolor{nred}{Intr. à Probab. e Estatística} &
    \textcolor{nblue}{Álgebra Linear} &
    \textcolor{nblue}{Análise de Algoritmos II} \\
    \textcolor{nred}{Funções de Uma Variável} &
    \textcolor{nblue}{Matemática Discreta} &
    \textcolor{nblue}{Ling. Formais e Autômatos} \\
    \textcolor{nred}{Funções de Várias Variáveis} &
    \textcolor{nblue}{Matemática Discreta II} &
    \textcolor{nblue}{Engenharia de Software}\\
    \textcolor{nblue}{Otimização Linear} & & \\
    \hline
    
    \multicolumn{3}{p{0.95\textwidth}}{}\\
    \multicolumn{3}{p{0.95\textwidth}}{Gerir a sua própria aprendizagem e
    desenvolvimento, incluindo a gestão de tempo e competências
    organizacionais}\\
    \hline
    \textcolor{nblue}{Trab. de Concl. de Curso I} &
    \textcolor{nblue}{Trab. de Concl. de Curso II} &
    \textcolor{nblue}{Trab. de Concl. de Curso III} \\
    \hline
    
    \multicolumn{3}{p{0.95\textwidth}}{}\\
    \multicolumn{3}{p{0.95\textwidth}}{Preparar e apresentar seus trabalhos e
    problemas técnicos e suas soluções para audiências diversas, em formatos
    apropriados (oral e escrito)}\\
    \hline
    \textcolor{nblue}{\small Metod. e Escr. Cient. para Comp.} &
    \textcolor{nblue}{Trab. de Concl. de Curso I} &
    \textcolor{nblue}{Trab. de Concl. de Curso II}\\
    \textcolor{nblue}{Engenharia de Software} &
    \textcolor{nblue}{Trab. de Concl. de Curso III} & \\
    \hline
    
    \multicolumn{3}{p{0.95\textwidth}}{}\\
    \multicolumn{3}{p{0.95\textwidth}}{Avaliar criticamente projetos de sistemas de computação}\\
    \hline
    \textcolor{nred}{Ciência, Tecnologia e Sociedade} &
    \textcolor{nblue}{Segurança de Dados} &
    \textcolor{nblue}{Análise de Algoritmos I} \\
    \textcolor{nblue}{Engenharia de Software} &
    \textcolor{nblue}{Comput., Ética e Sociedade} &
    \textcolor{nblue}{Análise de Algoritmos II} \\
    \textcolor{nblue}{Sistemas Distribuídos} & 
    \textcolor{nblue}{Redes de Computadores} & \\
    \hline
    
    \multicolumn{3}{p{0.95\textwidth}}{}\\
    \multicolumn{3}{p{0.95\textwidth}}{Adequar-se rapidamente às mudanças
    tecnológicas e aos novos ambientes de trabalho}\\
    \hline
    \textcolor{nblue}{Comput., Ética e Sociedade} & & \\
    \hline
    
    \multicolumn{3}{p{0.95\textwidth}}{}\\
    \multicolumn{3}{p{0.95\textwidth}}{Ler textos técnicos na língua inglesa}\\
    \hline
    \textcolor{nblue}{\small Metod. e Escr. Cient. para Comp.} &
    \textcolor{nblue}{Trab. de Concl. de Curso I} &
    \textcolor{nblue}{Trab. de Concl. de Curso II} \\
    \textcolor{nblue}{Trab. de Concl. de Curso III} & & \\
    \hline
    
    \multicolumn{3}{p{0.95\textwidth}}{}\\
    \multicolumn{3}{p{0.95\textwidth}}{Empreender e exercer liderança, coordenação
    e supervisão na sua área de atuação profissional}\\
    \hline
    \textcolor{npurple}{Ações de extensão} && \\
    \hline
    
    \multicolumn{3}{p{0.95\textwidth}}{}\\
    \multicolumn{3}{p{0.95\textwidth}}{Ser capaz de realizar trabalho cooperativo e entender os
    benefícios que este pode produzir}\\
    \hline
    \textcolor{nred}{Ciência, Tecnologia e Sociedade} &
    \textcolor{nblue}{Comput., Ética e Sociedade} &
    \textcolor{nblue}{Engenharia de Software}\\
    \hline
    
    \multicolumn{3}{p{0.95\textwidth}}{}\\
    \multicolumn{3}{p{0.95\textwidth}}{Compreender os fatos essenciais, os
    conceitos, os princípios e as teorias relacionadas à Ciência da Computação
    para o desenvolvimento de software e hardware e suas aplicações}\\
    \hline
    \textcolor{nred}{Bases Comput. da Ciência} &
    \textcolor{nblue}{Progr. Funcional} & 
    \textcolor{nblue}{Algs. em Grafos} \\
    \textcolor{nred}{Processamento da Informação} &
    \textcolor{nblue}{Algs. e Estruturas de Dados I} &
    \textcolor{nblue}{Otimização Linear}\\
    \textcolor{nblue}{Progr. Estruturada} & 
    \textcolor{nblue}{Algs. Estruturas de Dados II} &
    \textcolor{nblue}{Circuitos Digitais}\\
    \textcolor{nblue}{Análise de Algoritmos I} &
    \textcolor{nblue}{Ling. Formais e Autômatos} &
    \textcolor{nblue}{Sistemas Digitais}\\
    \textcolor{nblue}{Análise de Algoritmos II} &
    \textcolor{nblue}{Matemática Discreta} &
    \textcolor{nblue}{Sistemas Operacionais}\\
    \textcolor{nblue}{Arquitetura de Computadores} & 
    \textcolor{nblue}{Matemática Discreta II} & \\
    \hline
    
    \multicolumn{3}{p{0.95\textwidth}}{}\\
    \multicolumn{3}{p{0.95\textwidth}}{Reconhecer a importância do pensamento
    computacional no cotidiano e sua aplicação em circunstâncias apropriadas e
    em domínios diversos}\\
    \hline
    \textcolor{nred}{Comunicação e Redes} &
    \textcolor{nblue}{Progr. Estruturada} &
    \textcolor{nblue}{Matemática Discreta} \\
    \textcolor{nred}{Processamento da Informação} &
    \textcolor{nblue}{Algs. e Estruturas de Dados I}&
    \textcolor{nblue}{Matemática Discreta II}\\
    \textcolor{nred}{Ciência, Tecnologia e Sociedade} &
    \textcolor{nblue}{Algs. e Estruturas de Dados II}&
    \textcolor{nblue}{Algs. em Grafos}\\
    \textcolor{nred}{Bases Comput. da Ciência} &
    \textcolor{nblue}{Comput., Ética e Sociedade} &
    \textcolor{nblue}{Otimização Linear}\\
    \hline
    
    \multicolumn{3}{p{0.95\textwidth}}{}\\
    \multicolumn{3}{p{0.95\textwidth}}{Identificar e gerenciar os riscos que podem
    estar envolvidos na operação de equipamentos de computação (incluindo os
    aspectos de dependabilidade e segurança)}\\
    \hline
    \textcolor{nblue}{Segurança de Dados} &
    \textcolor{nblue}{Comput., Ética e Sociedade} &
    \textcolor{nblue}{Circuitos Digitais}\\
    \textcolor{nblue}{\small Model. de Banco de Dados} &
    \textcolor{nblue}{Redes de Computadores} & \\
    \hline
    
    \multicolumn{3}{p{0.95\textwidth}}{}\\
    \multicolumn{3}{p{0.95\textwidth}}{Identificar e analisar requisitos e
    especificações para problemas específicos e planejar estratégias para suas
    soluções}\\
    \hline
    \textcolor{nblue}{Engenharia de Software} &
    \textcolor{nblue}{Arquitetura de Computadores} &
    \textcolor{nblue}{Circuitos Digitais} \\
    \textcolor{nblue}{Análise de Algoritmos I} &
    \textcolor{nblue}{\small Model. de Banco de Dados} &
    \textcolor{nblue}{Sistemas Digitais} \\
    \textcolor{nblue}{Análise de Algoritmos II} & & \\
    \hline
    
    \multicolumn{3}{p{0.95\textwidth}}{}\\
    \multicolumn{3}{p{0.95\textwidth}}{Especificar, projetar, implementar, manter e
    avaliar sistemas de computação, empregando teorias, práticas e ferramentas
    adequadas}\\
    \hline
    \textcolor{nblue}{Sistemas Operacionais} &
    \textcolor{nblue}{Arquitetura de Computadores} &
    \textcolor{nblue}{Análise de Algoritmos I} \\
    \textcolor{nblue}{\small Model. de Banco de Dados} &
    \textcolor{nblue}{Redes de Computadores} &
    \textcolor{nblue}{Análise de Algoritmos II} \\
    \textcolor{nblue}{Engenharia de Software} &
    \textcolor{nblue}{Sistemas Distribuídos} &
    \textcolor{nblue}{Sistemas Digitais}\\
    \textcolor{nblue}{Compiladores e Interpretadores} & & \\
    \hline
    
    \multicolumn{3}{p{0.95\textwidth}}{}\\
    \multicolumn{3}{p{0.95\textwidth}}{Conceber soluções computacionais a partir de
    decisões visando o equilíbrio de todos os fatores envolvidos}\\
    \hline
    \textcolor{nred}{Processamento da Informação} &
    \textcolor{nblue}{Algs. e Estruturas de Dados I}&
    \textcolor{nblue}{Segurança de Dados}\\
    \textcolor{nred}{Ciência, Tecnologia e Sociedade} &
    \textcolor{nblue}{Algs. e Estruturas de Dados II}&
    \textcolor{nblue}{Engenharia de Software}\\
    \textcolor{nblue}{Progr. Estruturada} &
    \textcolor{nblue}{Comput., Ética e Sociedade}&
    \textcolor{nblue}{Prog. Funcional} \\
    \textcolor{nblue}{Inteligência Artificial} &
    \textcolor{nblue}{Compiladores e Interpretadores} & \\
    \hline
    
    \multicolumn{3}{p{0.95\textwidth}}{}\\
    \multicolumn{3}{p{0.95\textwidth}}{Empregar metodologias que visem garantir
    critérios de qualidade ao longo de todas as etapas de desenvolvimento de
    uma solução computacional}\\
    \hline
    \textcolor{nblue}{Engenharia de Software} & & \\
    \hline
    
    \multicolumn{3}{p{0.95\textwidth}}{}\\
    \multicolumn{3}{p{0.95\textwidth}}{Analisar quanto um sistema baseado em
    computadores atende os critérios definidos para seu uso corrente e futuro
    (adequabilidade)}\\
    \hline
    \textcolor{nblue}{Engenharia de Software} &
    \textcolor{nblue}{Análise de Algoritmos I} &
    \textcolor{nblue}{\small Model. de Banco de Dados}\\
    \hline
    
    \multicolumn{3}{p{0.95\textwidth}}{}\\
    \multicolumn{3}{p{0.95\textwidth}}{Gerenciar projetos de desenvolvimento de
    sistemas computacionais}\\
    \hline
    \textcolor{nblue}{Compiladores e Interpretadores} &
    \textcolor{nblue}{Sistemas Operacionais}&
    \textcolor{nblue}{\small Model. de Banco de Dados}\\
    \textcolor{nblue}{Engenharia de Software} &
    \textcolor{nyellow}{Gestão de Projetos de Software} & \\
    \hline
    
    \multicolumn{3}{p{0.95\textwidth}}{}\\
    \multicolumn{3}{p{0.95\textwidth}}{Aplicar temas e princípios recorrentes, como
    abstração, complexidade, princípio de localidade de referência (caching),
    compartilhamento de recursos, segurança, concorrência, evolução de
    sistemas, entre outros, e reconhecer que esses temas e princípios são
    fundamentais à área de Ciência da Computação}\\
    \hline
    \textcolor{nblue}{Análise de Algoritmos I} &
    \textcolor{nblue}{Progr. Orientada a Objetos} &
    \textcolor{nblue}{\small Model. de Banco de Dados}\\
    \textcolor{nblue}{Análise de Algoritmos II} &
    \textcolor{nblue}{Segurança de Dados} &
    \textcolor{nblue}{Sistemas Distribuídos}\\
    \textcolor{nblue}{Redes de Computadores} & 
    \textcolor{nblue}{Arquitetura de Computadores} &
    \textcolor{nblue}{Sistemas Operacionais}\\
    \textcolor{nblue}{Inteligência Artificial} & 
    \textcolor{nblue}{Engenharia de Software} &
    \textcolor{nblue}{Sistemas Digitais}\\
    \hline
    
    
    \multicolumn{3}{p{0.95\textwidth}}{}\\
    \multicolumn{3}{p{0.95\textwidth}}{Escolher e aplicar boas práticas e técnicas
    que conduzam ao raciocínio rigoroso no planejamento, na execução e no
    acompanhamento, na medição e gerenciamento geral da qualidade de sistemas
    computacionais}\\
    \hline
    \textcolor{nblue}{Análise de Algoritmos I} &
    \textcolor{nblue}{Comput., Ética e Sociedade} &
    \textcolor{nblue}{\small Model. de Banco de Dados}\\
    \textcolor{nblue}{Análise de Algoritmos II} &
    \textcolor{nyellow}{Sistemas de Informação} &
    \textcolor{nblue}{Sistemas Digitais}\\
    \textcolor{nblue}{Engenharia de Software} &
    \textcolor{nblue}{Compiladores e Interpretadores} & \\
    \hline
    
    
    \multicolumn{3}{p{0.95\textwidth}}{}\\
    \multicolumn{3}{p{0.95\textwidth}}{Aplicar os princípios de gerência,
    organização e recuperação da informação de vários tipos, incluindo texto,
    imagem, som e vídeo}\\
    \hline
    \textcolor{nblue}{\small Model. de Banco de Dados} &
    \textcolor{nblue}{Algs. e Estruturas de Dados I} & 
    \textcolor{nblue}{Sistemas Distribuídos}\\
    \textcolor{nblue}{Computação Gráfica} &
    \textcolor{nblue}{Algs. e Estruturas de Dados II} &
    \textcolor{nyellow}{Proc. de Sinais Neurais} \\
    \textcolor{nblue}{Redes de Computadores} & 
    \textcolor{nyellow}{Process. Digital de Imagens} & \\
    \hline
    
    \multicolumn{3}{p{0.95\textwidth}}{}\\
    \multicolumn{3}{p{0.95\textwidth}}{Aplicar os princípios de interação
    humano-computador para avaliar e construir uma grande variedade de
    produtos, incluindo interface de usuário, páginas WEB, sistemas multimídia
    e sistemas móveis}\\
    \hline
    \textcolor{nblue}{Computação Gráfica} & 
    \textcolor{nyellow}{\small Interação Humano-Computador} &
    \textcolor{nyellow}{Sistemas Multimidia}\\
    \textcolor{nyellow}{Programação para Web} &
    \textcolor{nyellow}{\small Prog. Av. de Dispositivos Móveis} & 
    \textcolor{nyellow}{Visão Computacional}\\
    \textcolor{nyellow}{Sistemas Inteligentes} & & \\
    \hline
    
\end{longtable}


\newpage
\section{Ações acadêmicas complementares à formação}
\label{sec:acoes_complementares}

A UFABC possui diversos programas e ações para promover a qualidade do ensino
de graduação, dentre os quais citamos:
\begin{description}
    \item[Projeto de Ensino-Aprendizagem Tutorial -- PEAT:] tem como objetivo
    promover a adaptação do aluno ao projeto acadêmico da UFABC, orientando-o
    para uma transição tranquila e organizada do Ensino Médio para o Superior.
    Mais informações em: \url{http://prograd.ufabc.edu.br/peat};

    \item[Programas de iniciação científica:] têm como objetivo introduzir os
    alunos de graduação na pesquisa científica, visando fundamentalmente
    colocar o aluno desde cedo em contato direto com a atividade científica e
    engajá-lo na pesquisa. A UFABC possui as seguintes modalidades:
    \begin{itemize}
        \item Programa Pesquisando Desde o Primeiro Dia -- PDPD;
        \item Programa de Iniciação Científica -- PIC/UFABC;
        \item Programa Institucional de Bolsas de Iniciação Científica -- PIBIC/CNPq;
        \item Programa Institucional de Bolsas de Iniciação Científica -- PIBIC/CNPq nas Ações Afirmativas.
    \end{itemize}
    Mais informações em:
    \url{https://propes.ufabc.edu.br/perfis-de-acesso/aluno};

    \item[Programa de monitoria acadêmica:] têm como objetivo selecionar alunos
    para desenvolverem atividades de monitoria.
    Mais informações em:
    \url{http://prograd.ufabc.edu.br/monitoria-academica};

    \item[Programa Institucional de Bolsas de Iniciação à Docência -- PIBID:] é
    um programa da Coordenação de Aperfeiçoamento de Pessoal de Nível Superior
    (CAPES) que tem por finalidade fomentar a iniciação à docência,
    contribuindo para o aperfeiçoamento da formação de docentes em nível
    superior e para a melhoria da qualidade da educação básica pública
    brasileira.
    Mais informações em: \url{https://pibid.ufabc.edu.br/};

    \item[Ações extensionistas:] esse tipo de atividade ultrapassa o âmbito
    específico de atuação do Instituto no que se refere ao Ensino (Graduação e
    Pós-Graduação) e Pesquisa. A Extensão é uma das funções sociais da
    Universidade, realizada por meio de um conjunto de ações dirigidas à
    sociedade, as quais devem estar indissociavelmente vinculadas ao Ensino e à
    Pesquisa.
    Mais informações em: \url{http://proec.ufabc.edu.br/};

    \item[Programa de Educação Tutorial -- PET:] tem como proposta desenvolver
    atividades que propiciem a ciência, tecnologia e inovação de dentro para
    fora da Universidade, conscientizando seus discentes da sua importância e
    de como fazer, assim como proporcionar ao corpo docente um ambiente
    favorável ao seu desenvolvimento e dar acesso a qualquer comunidade a esse
    recurso tanto acadêmica quanto externamente.
    Mais informações em:
    \url{http://prograd.ufabc.edu.br/pet};

    \item[Cursos de língua estrangeira:] oferecidos pelo Núcleo Educacional de
    Tecnologias e Línguas. Mais informações em:
    \url{https://netel.ufabc.edu.br/};

    \item[Mobilidade acadêmica:] consiste em um período de estudos, em regra de
    1 semestre, em uma universidade estrangeira ou nacional, com o objetivo de
    oferecer ao aluno experiências enriquecedoras capazes de agregar
    positivamente sua vida acadêmica, profissional e pessoal.
    Mais informações em: \url{https://ri.ufabc.edu.br/mobilidade-academica};

    \item[Monitoria inclusiva:] os monitores inclusivos são alunos de graduação
    que se dedicam 10 horas semanais em atividades de ações afirmativas ao
    aluno com deficiência, dando suporte como ledor, transcritor,
    audiodescritora de figuras, imagens, desenhos e vídeos. Outra atividade que
    também demanda atenção do monitor inclusivo é a adaptação de materiais e
    livros usados por alunos com deficiência visual.
    Mais informações em:
    \url{https://proap.ufabc.edu.br/acessibilidade-ufabc/servicos-e-recursos/monitoria-inclusiva};

    \item[Programa de Apoio ao Desenvolvimento Acadêmico -- PADA:] realiza
    atividades de orientação pedagógica a discentes de graduação nas áreas de:
    planejamento dos estudos junto a estudantes dos Bacharelados
    Interdisciplinares (BIs) e Licenciaturas Interdisciplinares (LIs);
    requisitos para integralização dos BIs e LIs; prazos para conclusão dos
    cursos interdisciplinares; prevenção ao desligamento dos cursos
    interdisciplinares.
    Mais informações em: \url{https://prograd.ufabc.edu.br/pada}.
\end{description}


\newpage
\section {Atividades Complementares}

As atividades complementares são todas as atividades de diversas naturezas, que não se incluem no desenvolvimento regular das disciplinas constantes na matriz curricular do BCC, mas que são relevantes para a formação do aluno.

O objetivo do incentivo à realização de atividades complementares consiste em fornecer ao estudante a oportunidade de enriquecer sua formação com a participação em atividades de natureza diversificada. Como consequência, tem-se a acentuação do caráter interdisciplinar e amplo da formação do aluno, além do fortalecimento do vínculo entre teoria e prática.

Uma vez que o BC\&T é requisito para o BCC, e neste curso já está prevista a realização de 120 horas de atividades complementares, o BCC não exigirá a realização de atividades complementares específicas além das já previstas no BC\&T. 

A forma de validação da carga horária dessas atividades encontra-se na Resolução CG no. 11, de 28 de junho de 2016, publicado pelo Boletim de Serviços no. 568, de 05 de julho de 2016.


\newpage
\section{Estágio curricular}

O Estágio Curricular não é obrigatório para o curso de Bacharelado em Ciência
da Computação.


\newpage
\section {Trabalho de Conclusão de Curso}


\newpage
\section{Avaliação de processo ensino-aprendizagem}
\label{sec:avaliacao_ensino_aprendizagem}

A avaliação do processo de ensino e aprendizagem dos discentes na UFABC é feita
por meio de conceitos, o que permite uma análise mais qualitativa do
aproveitamento do aluno.
Segundo a Resolução ConsEPE nº 147, de 19 de março de 2013, os coeficientes de
desempenho utilizados na Instituição consistem em:
\begin{description}
    \item[A -] Desempenho excepcional, demonstrando excelente compreensão da
    disciplina e do uso do conteúdo.
    \item[B -] Bom desempenho, demonstrando boa capacidade de uso dos conceitos
    da disciplina.
    \item[C -] Desempenho mínimo satisfatório, demonstrando capacidade de uso
    adequado dos conceitos da disciplina, habilidade para enfrentar problemas
    relativamente simples e prosseguir em estudos avançados.
    \item[D -] Aproveitamento mínimo não satisfatório dos conceitos da
    disciplina, com familiaridade parcial do assunto e alguma capacidade para
    resolver problemas simples, mas demonstrando deficiências que exigem
    trabalho adicional para prosseguir em estudos avançados. Nesse caso, o
    aluno é aprovado na expectativa de que obtenha um conceito melhor em outra
    disciplina, para compensar o conceito D no cálculo do CR. Havendo vaga, o
    aluno poderá cursar esta disciplina novamente.
    \item[F -] Reprovado. A disciplina deve ser cursada novamente para obtenção
    de crédito.
    \item[O -] Reprovado por falta. A disciplina deve ser cursada novamente
    para obtenção de crédito.
\end{description}

Os conceitos a serem atribuídos aos estudantes, em uma dada disciplina, não
deverão estar rigidamente relacionados a qualquer nota numérica de provas,
trabalhos ou exercícios.
Os resultados também considerarão a capacidade do aluno de utilizar os
conceitos e material das disciplinas, criatividade, originalidade, clareza de
apresentação e participação em sala de aula e/ou laboratórios. 
O aluno, ao iniciar uma disciplina, será informado sobre as normas e critérios
de avaliação que serão considerados.

Não há um limite mínimo de avaliações a serem realizadas, mas, dado o caráter
qualitativo do sistema, é indicado que sejam realizadas ao menos duas em cada
disciplina durante o período letivo.
Esse mínimo sugere a possibilidade de ser feita uma avaliação diagnóstica logo
no início do período letivo, que identifique a capacidade do aluno em lidar com
conceitos que apoiarão o desenvolvimento de novos conhecimentos e o quanto ele
conhece dos conteúdos a serem discutidos na disciplina, e outra no final do
período, que possa identificar a evolução do aluno relativamente ao estágio de
diagnóstico inicial. 
De posse do diagnóstico inicial, o próprio professor poderá ser mais eficiente
na mediação com os alunos no desenvolvimento da disciplina. 

Serão apoiadas e incentivadas as iniciativas de se gerar novos documentos de
avaliação, como atividades extraclasse, tarefas em grupo, listas de exercícios,
atividades em sala e/ou em laboratório, observações do professor,
auto-avaliação, seminários, exposições, projetos, sempre no intuito de se
viabilizar um processo de avaliação que não seja apenas qualitativo, mas que se
aproxime de uma avaliação contínua.
Assim, propõe-se não apenas a avaliação de conteúdos, mas de estratégias
cognitivas e habilidades e competências desenvolvidas. 

Por fim, deverá ser levada em alta consideração o processo evolutivo descrito
pelas sucessivas avaliações no desempenho do aluno para que se faça a
atribuição de um conceito a ele.

Cabe ressaltar que os critérios de recuperação do curso da UFABC são atualmente
regulamentados pela Resolução ConsEPE nº 182 (ou outra resolução que venha a
substituí-la).





\newpage
\section{Infraestrutura}

A UFABC é uma universidade multicampi.
Tanto o campus de Santo André como o campus de São Bernardo do Campo possuem
biblioteca, conexão de internet de alta velocidade, laboratórios didáticos de
experimentação, de ensino e computação, laboratórios de pesquisa, biotérios de
criação e manutenção de animais de experimentação, setores administrativos,
salas de reunião e salas de docentes.

\subsection{Campus Santo André}

O `Bloco A' de edifícios do Campus Santo André mede cerca de 39.000 m$^2$ onde
está localizada a maior parte das salas de aula, laboratórios de pesquisa e
salas de docentes daquele campus.
Esta obra possui três torres principais, cada um relacionado a um centro desta
universidade: Centro de Engenharias, Modelagem e Ciências Sociais Aplicadas
(CECS), Centro de Ciências Naturais e Humanas (CCNH) e Centro de Matemática,
Computação e Cognição (CMCC). 
As três edificações estão interligadas por áreas comuns nos primeiros três
andares de cada prédio.
Nessas áreas comuns estão instaladas salas de aula da graduação e setores
administrativos.
A ideia de continuidade física entre as áreas da UFABC está em consonância com
seu projeto de criação que visa a interdisciplinaridade como sua principal
meta.
Algumas salas de docentes, laboratórios didáticos e de pesquisa, e salas de
aula também estão localizados no prédio de 11 andares adjacente ao `Bloco A',
denominado `Bloco B'.
Por fim, O `Bloco L', com área construída de mais de 16.800 m$^2$ abriga 72
laboratórios didáticos e de pesquisa dos três Centros, além de lanchonetes,
almoxarifado entre outros.

\subsection{Campus São Bernardo do Campo} 

O campus de São Bernardo do Campo possui laboratórios didáticos para
experimentação e computação nos Blocos `Alfa' e `Tau'.
O Bloco `Beta' abriga a biblioteca, anfiteatros e um amplo auditório de 400
lugares.
Estão alocados nos laboratórios didáticos do bloco Alfa diversos modelos
anatômicos e sistemas de ensino de fisiologia (i-Works).
Estão previstos ainda outros edifícios, já em construção, abrigando
laboratórios didáticos específicos das Engenharias (Bloco `Omega'),
laboratórios de pesquisa (Bloco `Zeta') e um biotério de caráter multiusuário
de criação e manutenção de animais de experimentação.

\subsection{Laboratórios didáticos}

A Pró-Reitoria de Graduação possui em sua infraestrutura um grupo de
laboratórios compartilhados por todos os cursos de graduação. 
A Coordenadoria dos Laboratórios Didáticos (CLD), vinculada à PROGRAD, é
responsável pela gestão administrativa dos laboratórios didáticos e por
realizar a interface entre docentes, discentes e técnicos de laboratório nas
diferentes áreas, de forma a garantir o bom andamento dos cursos de graduação,
no que se refere às atividades práticas em laboratório.

A CLD é composta por um Coordenador dos Laboratórios Úmidos, um Coordenador dos
Laboratórios Secos e um Coordenador dos Laboratórios de Informática e Práticas
de Ensino, bem como equipe técnico-administrativa. 

Dentre as atividades da CLD destacam-se o atendimento diário a toda comunidade
acadêmica; a elaboração de Política de Uso e Segurança dos Laboratórios
Didáticos e a análise e adequação da alocação de turmas nos laboratórios em
cada quadrimestre letivo, garantindo a adequação dos espaços às atividades
propostas em cada disciplina e melhor utilização de recursos da UFABC, o
gerenciamento da infraestrutura dos laboratórios didáticos, materiais, recursos
humanos, treinamento, encaminhamento para manutenção preventiva e corretiva de
todos os equipamentos. 

Os laboratórios são dedicados às atividades didáticas práticas que necessitam
de infraestrutura específica e diferenciada, não atendidas por uma sala de aula
convencional. 
São quatro diferentes categorias de laboratórios didáticos disponíveis para os
usos dos cursos de graduação da UFABC: 
\begin{itemize}
    \item \underline{Laboratórios Didáticos Secos} são espaços destinados às
    aulas da graduação que necessitem de uma infraestrutura com bancadas e
    instalação elétrica e/ou instalação hidráulica e/ou gases, uso de kits
    didáticos e mapas, entre outros.
    \item \underline{Laboratórios Didáticos Úmidos} são espaços destinados às
    aulas da graduação que necessitem manipulação de agentes químicos ou
    biológicos, uma infraestrutura com bancadas de granito, com capelas de
    exaustão e com instalações hidráulica, elétrica e de gases
    \item \underline{Laboratórios Didáticos Práticas de Ensino} são espaços
    destinados ao suporte dos cursos de licenciatura, desenvolvimento de
    habilidades e competências para docência da educação básica, podendo ser
    úteis também para desenvolvimentos das habilidades e competências para
    docência do ensino superior.
    \item \underline{Laboratórios Didáticos de Informática} são espaços para
    aulas utilizando recursos de tecnologia de informação como
    microcomputadores, acesso à internet, linguagens de programação, softwares,
    hardwares e periféricos.
\end{itemize}

Anexo a cada laboratório há uma sala de suporte técnico que acomoda quatro técnicos
de laboratório, cumprindo as seguintes funções:
\begin{itemize}
    \item nos períodos extra aula, auxiliam a comunidade no que diz respeito a
    atividades de graduação, pós-graduação e extensão em suas atividades
    práticas (projetos de disciplinas, iniciação científica, mestrado e
    doutorado), participam dos processos de compras levantando a minuta dos
    materiais necessários, fazem controle de estoque de materiais, bem como
    cooperam com os professores durante a realização de testes e experimentos
    que serão incorporados nas disciplinas e preparação do laboratório para a
    aula prática;
    \item nos períodos de aula, oferecem apoio para os professores e alunos
    durante o experimento, repondo materiais, auxiliando no uso de equipamentos
    e prezando pelo bom uso dos materiais de laboratório. Para isso, os
    técnicos são alocados previamente em determinadas disciplinas, conforme a
    sua formação (eletrônica, eletrotécnica, materiais, mecânica, mecatrônica,
    edificações, química, biologia, informática, etc).
\end{itemize}
Os técnicos trabalham em esquema de horários alternados, possibilitando o apoio
às atividades práticas ao longo de todo período de funcionamento da UFABC, das
8h às 23h.

Além dos técnicos, a sala de suporte armazena alguns equipamentos e kits
didáticos utilizados nas disciplinas. 
Há também a sala de suporte técnico, que funciona como almoxarifado,
armazenando demais equipamentos e kits didáticos utilizados durante o
quadrimestre.

A UFABC dispõe ainda de uma oficina mecânica de apoio, com quatro técnicos
especializados na área, e atende a demanda de toda comunidade acadêmica
(centros, graduação, extensão e prefeitura universitária), para a construção e
pequenas reparações de kits didáticos e dispositivos para equipamentos usados
na graduação e pesquisa, além do auxílio a discentes na construção e montagem de
trabalhos de graduação e pós, projetos acadêmicos como BAJA, Aerodesign, etc.
A oficina mecânica atende no horário das 8h às 17h.
Esta oficina está equipada com as seguintes máquinas operatrizes: torno CNC,
centro de usinagem CNC, torno mecânico horizontal, fresadora universal,
retificadora plana, furadeira de coluna, furadeira de bancada, esmeril, serra
de fita vertical, lixadeira, serra de fita horizontal, prensa hidráulica,
máquina de solda elétrica TIG, aparelho de solda oxi-acetilênica, calandra,
curvadora de tubos, guilhotina e dobradora de chapas. 
Além disso, a oficina mecânica possui duas bancadas e uma grande variedade de
ferramentas para trabalhos manuais: chaves para aperto, limas, serras manuais,
alicates de diversos tipos, torquímetros, martelos e diversas ferramentas de
corte de uso comum em mecânica, como também, ferramentas manuais elétricas como
furadeiras manuais, serra tico-tico, grampeadeira, etc. 
Também estão disponíveis vários tipos de instrumentos de medição comuns em
metrologia: paquímetros analógicos e digitais, micrômetros analógicos com
batentes intercambiáveis, micrômetros para medição interna, esquadros e
goniômetros, traçadores de altura, desempeno, escalas metálicas, relógios
comparadores analógicos e digitais e calibradores.
Com estes equipamentos e ferramentas, é possível a realização de uma ampla gama
de trabalhos de usinagem, ajustes, montagem e desmontagem de máquinas e
componentes mecânicos.

A alocação de laboratórios didáticos para as turmas de disciplinas com carga
horária prática ou aquelas que necessitem do uso de um laboratório é feita pelo
coordenador do curso, a cada quadrimestre, durante o período estipulado pela
Pró-Reitoria de Graduação.

O docente da disciplina com carga horária alocada nos laboratórios didáticos é
responsável pelas aulas práticas da disciplina, não podendo se ausentar do
laboratório durante a aula prática.
Atividades como treinamentos, instalação ou manutenção de equipamentos nos
laboratórios didáticos ou aulas pontuais são previamente agendadas com a equipe
técnica responsável e acompanhadas por um técnico de laboratório.

Como os laboratórios são compartilhados, todos os cursos podem realizar de
diferentes atividades didáticas dentro dos diversos laboratórios, otimizando o
uso dos recursos materiais e ampliando as possibilidades didáticas dos docentes
da UFABC e a prática da interdisciplinaridade, respeitando as necessidades de
cada disciplina ou aula de acordo com a classificação do laboratório e dos
materiais e equipamentos disponíveis nele.

\subsection{Sistema de Bibliotecas – SISBI}

O Sistema de Bibliotecas da UFABC, cuja finalidade é atender as demandas
informacionais da comunidade universitária e científica interna e externa à
Universidade, é formado por unidades de bibliotecas localizadas nos Campi de
Santo André e São Bernardo do Campo, responsáveis por atender e apoiar a
comunidade universitária em suas atividades de ensino pesquisa e extensão, de
forma articulada e pautada na proposta interdisciplinar do projeto pedagógico e
de seu plano de desenvolvimento institucional.

As Bibliotecas que compõem o Sistema possuem amplo e diversificado acervo, com
aproximadamente 100.000 exemplares de livros físicos e 42.000 títulos de livros
eletrônicos, sendo todas as coleções da editora Springer Nature entre os anos
de 2005 e 2014, todos os títulos publicados pela editora Wiley em 2016 e
pelos títulos da editora Ebsco referentes a coleção EbscoHost. 
Além disso, conta com títulos resultantes de assinaturas anuais com demais
editoras, como Elsevier, Cengage Learning e Wiley, além de ter uma filmoteca,
que conta com mais de 1000 títulos de filmes.
O SisBi ainda dispõe do sistema SophiA, que permite o acesso ao seu catálogo e
portal na internet para acesso às informações sobre seus serviços e a conteúdos
externos, como: sistema Scielo, que contempla seleção de periódicos científicos
brasileiros; sistema Biblioteca Digital Brasileira de Teses e Dissertações
(BDTD); sistema COMUT, que permite a obtenção de cópias de documentos
técnico-científicos disponíveis nos acervos das principais bibliotecas
brasileiras e em serviços de informações internacionais; Portal de Periódicos
da CAPES, que oferece uma seleção das mais importantes fontes de informação
científica e tecnológica, de acesso gratuito na Web.
Atualmente, o portal dispõe de 34.457 periódicos eletrônicos, relacionados às
diversas áreas do conhecimento e, ainda, acesso a mais de 2.000 bases de dados;
dentre outros.

Convênios também são estabelecidos pelo SisBi, entre os mais significativos o
serviço de Empréstimo Entre Bibliotecas (EEB), que estabelece a cooperação e
potencializa a utilização do acervo das instituições universitárias
participantes, favorecendo a disseminação da informação entre universitários e
pesquisadores de todo o país. Outro convênio a ser notado é com o IBGE, que tem
por objetivo ampliar para a sociedade, o acesso às informações produzidas por
meio de cooperação técnica com o Centro de Documentação e Disseminação de
Informações do IBGE.
Assim, o SisBi passou a ser depositário das publicações editadas por esse órgão.

As unidades de bibliotecas atendem a comunidade de segunda a sexta, das 8h às
22h, mantendo-se em uma estrutura física com área total de 4.529 m², onde se
distribuem 521 assentos, além de terminais de consulta ao acervo.
Buscando promover o exercício a reflexão crítica nos espaços universitários,
bem como a interação com os diversos públicos, desenvolve ainda programas e
projetos culturais como CineArte, exibido também ao ar livre, PublicArte,
Saraus e Exposições.

\subsection{Tecnologias digitais}

As Tecnologias de Informação e Comunicação (TIC) têm sido cada vez mais
utilizadas no processo de ensino e aprendizagem. Sua importância não está
restrita apenas à oferta de disciplinas e cursos semipresenciais, ou totalmente
a distância, ocupando um espaço importante também como mediadoras em
disciplinas e cursos presenciais.
As salas de aula da UFABC são equipadas com projetor multimídia e um
computador, e as disciplinas práticas, que demandam o uso de computadores e
internet, são ministradas em laboratórios equipados com 30-48 computadores com
acesso à Internet, projetor multimídia e softwares relacionados às atividades
desenvolvidas. 
Estão disponíveis também 10 lousas digitais, distribuídas em salas específicas
de cada centro. 
Para o uso dessas ferramentas e infraestrutura, os docentes contam com o
suporte técnico do Núcleo de Tecnologia da Informação (NTI) e da Coordenação de
Laboratórios Didáticos (CLD).

\subsection{Ambiente Virtual de Aprendizagem (AVA)}

Com o intuito de estimular a integração das TIC, a UFABC incentiva o uso de um
Ambiente Virtual de Aprendizagem (TIDIA 4 ou Moodle) como ferramenta de apoio
ao ensino presencial e semipresencial nas diversas disciplinas. 
O AVA pode possibilitar a interação entre alunos e professores por meio de
ferramentas de comunicação síncrona (e.g. bate papo/ chat) e assíncrona (e.g.
fórum de discussões, correio eletrônico), além de funcionar como repositório de
conteúdo didáticos, e permitir propostas de atividades individuais e
colaborativas.

\subsection{Núcleo Educacional de Tecnologias e Línguas (NETEL)}

No âmbito da utilização das TIC nas diferentes modalidades de ensino e
aprendizagem (presencial, semipresencial e a distância), a UFABC conta com o
apoio do Núcleo Educacional de Tecnologias e Línguas2
(\url{http://netel.ufabc.edu.br/}).
O NETEL está organizado em seis divisões (Cursos, Design e Inovação
Educacional, Tecnologias da Informação, Audiovisual, Comunicação e idiomas), e
oferece cursos de extensão e oficinas para capacitação de docentes interessados
na integração de novas metodologias e tecnologias digitais nas suas práticas de
ensino. 
Os cursos e oficinas são oferecidos periodicamente, nas modalidades presencial
e semipresencial, e possibilitam a formação e a atualização em diferentes
domínios, por exemplo: docência com tecnologias, desenvolvimento de objetos de
aprendizagem, jogos digitais educacionais, videoaulas, webconferência, lousa
digital, metodologias ativas de ensino, ferramentas digitais de apoio à
aprendizagem. 
Para apoiar a oferta de disciplinas na modalidade semipresencial, outras
iniciativas formativas do NETEL são os cursos semipresenciais Planejamento de
cursos virtuais, que se configura em uma oportunidade de reflexão e
compartilhamento de ideias sobre estratégias, ferramentas e métodos que apoiam
a criação de espaços virtuais de aprendizagem, e o curso Formação de Tutores
para EAD, que tem como objetivo capacitar alunos de graduação e pós-graduação e
pessoas interessadas em atuar como tutores/monitores. 
Para apoiar o docente na criação e oferta de disciplinas na modalidade
semipresencial, o NETEL conta com uma equipe de profissionais da área de Design
Instrucional e especialistas no desenvolvimento de recursos educacionais, como
objetos de aprendizagem e jogos educacionais. 
O NETEL possui também uma divisão de audiovisual com infraestrutura completa de
estúdio e equipamentos para gravação de videoaulas e podcasts. 
O estúdio proporciona apoio à comunidade acadêmica em diversos projetos de
extensão e outras iniciativas que demandam o uso de recursos audiovisuais como
filmagem de aulas, palestras etc. 
Em 2019 o NETEL passou a integrar uma nova divisão, divisão de idiomas, no qual
é responsável por desenvolver a política linguística da UFABC através de
ofertas de cursos de línguas gratuitos e presenciais como cursos de inglês,
português, espanhol e francês.

Por se tratar de uma instituição que busca excelência no uso das TIC, muitos
pesquisadores da UFABC têm desenvolvido pesquisas interdisciplinares nas áreas
de Educação, Ensino, Ciência da Computação, Comunicação etc., com o objetivo de
compreender as potencialidades de uso das TIC e sua influência nos processos de
ensino e aprendizagem. 
Neste contexto, os docentes envolvidos no núcleo juntamente com parceiros da
UFABC desenvolvem pesquisas com a finalidade de renovação e atualização
constante das TICs tanto no ensino quanto apoio ao mesmo.

\subsection{Oferta de disciplinas semipresenciais}

A Portaria MEC nº 2.117, de 6 de dezembro de 2019 (disponível em
\url{https://www.in.gov.br/en/web/dou/-/portaria-n-2.117-de-6-de-dezembro-de-2019-232670913}),
orienta sobre a oferta, por Instituições de Educação Superior (IES), de
disciplinas na modalidade à distância em cursos de graduação presencial. 
Neste sentido, as coordenações dos cursos de graduação juntamente com o seu
corpo docente poderão decidir como farão o uso desta portaria no sentido de
incluir componentes curriculares que, no todo ou em parte, utilizem a
modalidade de ensino semipresencial ou a distância, desde que esta oferta não
ultrapasse 40\% (quarenta por cento) da carga horária do curso. 
Uma mesma disciplina do curso poderá ser ofertada nos formatos presencial e
semipresencial, com Planos de Ensino devidamente adequados à sua oferta.
O número de créditos atribuídos a um componente curricular será o mesmo em
ambos os formatos. 
Para fins de registros escolares, não existe qualquer distinção entre as
ofertas presencial ou semipresencial de um dado componente curricular.
Portanto, em ambos os casos, as TICs, o papel dos tutores e dos docentes, a
metodologia de ensino, e o material didático a serem utilizados deverão ser
detalhados em proposta de Plano de Aula a ser avaliado pela coordenação do
curso antes de sua efetiva implantação. 
O uso desta portaria é de grande importância pois, motiva o uso das TICs nas
disciplinas de graduação favorecendo a renovação e modernização do ensino e
criando oportunidade para o desenvolvimento das habilidades digitais tanto dos
docentes quanto dos alunos da UFABC.

\subsection{Acessibilidade}

A UFABC possui um Núcleo de Acessibilidade, lotado na Pró-Reitoria de Assuntos
Comunitários e Políticas Afirmativas (ProAP) responsável por executar as
políticas de assistência estudantil direcionadas aos estudantes com deficiência
da nossa comunidade.
Essas ações e projetos visam eliminar as barreiras arquitetônicas, atitudinais
e de comunicação promovendo a inclusão das pessoas com deficiência. 
É papel da ProAP dar suporte a estudantes com deficiência ou necessidades
educacionais específicas, além de orientar a comunidade acadêmica nas questões
que envolvem o atendimento educacional destes estudantes. 
Além disso, a fim de possibilitar à pessoa com deficiência viver de forma
autônoma e participar de todos os aspectos da vida acadêmica, a ProAP preza
pela disseminação do conceito de desenho universal, conforme disposto na
legislação vigente. 
Orientar o corpo docente, acolher aos estudantes respeitando suas
especificidades, difundir e oferecer Tecnologias Assistivas, dar suporte de
monitoria acadêmica as disciplinas da graduação, disponibilizar tradução e
interpretação de LIBRAS, além da oferta de alguns programas de subsídios
financeiros propostos pelo Plano Nacional de Assistência Estudantil - PNAES,
também fazem parte dos programas em acessibilidade da UFABC. 


\newpage
\section {Docentes}

\begin{longtable}{|l|l|c|c|}
	\hline
	Docente & Área de formação & Titulação & dedicação\\
	\hline\hline
	Alexandre Donizeti Alves & Ciência da Computação & Doutor & DE \\
	Alexandre Noma & Ciência da Computação & Doutor & DE \\
	Ana Lígia Barbour Scott & Biofísica Molecular & Doutora & DE \\
	André Luiz Brandão & Ciência da Computação & Doutor & DE \\
	Antonio Sérgio Kimus Braz & Ciências Biológicas - Genética & Doutor & DE \\
	Aritanan Borges Garcia Gruber & Ciência da Computação - Pesquisa Operacional & Doutor & DE \\
	Bruno Augusto Dorta Marques & Ciência da Computação & Doutor & DE \\
	Carla Lopes Rodriguez & Artes Visuais & Doutora & DE \\
	Carla Negri Lintzmayer & Ciência da Computação & Doutora & DE \\
	Carlo Kleber da Silva Rodrigues & Engenharia de Sistemas e Computação & Doutor & DE \\
	Carlos Alberto Kamienski & Ciência da Computação & Doutor & DE \\
	Carlos da Silva dos Santos & Ciência da Computação & Doutor & DE \\
	Cláudio Nogueira de Meneses & Engenharia Industrial e de Sistemas & Doutor & DE \\
	Cristiane Maria Sato & Otimização e Combinatória & Doutora & DE \\
	Daniel Morgato Martin & Matemática & Doutor & DE \\
	David Corrêa Martins Júnior & Ciência da Computação & Doutor & DE \\
	Debora Maria Rossi de Medeiros & Ciência da Computação e Matemática Computacional & Doutora & DE \\
	Denis Gustavo Fantinato & Engenharia Elétrica & Doutor & DE \\
	Denise Hideko Goya & Ciência da Computação & Doutora & DE \\
	Diogo Santana Martins & Ciência da Computação e Matemática Computacional & Doutor & DE \\
	Edson Pinheiro Pimentel & Engenharia Eletrônica e Computação & Doutor & DE \\
	Emílio de Camargo Francesquini & Ciência da Computação & Doutor & DE \\
	Fabrício Olivetti de França & Engenharia Elétrica & Doutor & DE \\
	Fedor Pisnitchenko & Matemática Aplicada & Doutor & DE \\
	Fernando Teubl Ferreira & Engenharia Elétrica & Doutor & DE \\
	Flávio Eduardo Aoki Horita & Ciência da Computação e Matemática Computacional & Doutor & DE \\
	Francisco de Assis Zampirolli & Engenharia Elétrica & Doutor & DE \\
	Francisco Isidro Massetto & Engenharia Elétrica & Doutor & DE \\
	Francisco Javier Ropero Peláez & Neurociências & Doutor & DE \\
	Gordana Manic & Ciência da Computação & Doutora & DE \\
	Guiou Kobayashi & Engenharia Elétrica & Doutor & DE \\
	Gustavo Sousa Pavani & Engenharia Elétrica & Doutor & DE \\
	Harlen Costa Batagelo & Engenharia Elétrica & Doutor & DE \\
	Itana Stiubiener & Engenharia Elétrica  & Doutora & DE \\
	Jair Donadelli Júnior & Matemática Aplicada & Doutor & DE \\
	Jerônimo Cordoni Pellegrini & Ciência da Computação & Doutor & DE \\
	Jesús Pascual Mena-Chalco & Ciência da Computação & Doutor & DE \\
	João Marcelo Borovina Josko & Ciência da Computação & Doutor & DE \\
	João Paulo Gois & Ciência da Computação e Matemática Computacional & Doutor & DE \\
	João Ricardo Sato & Estatística & Doutor & DE \\
	José Artur Quilici Gonzalez & Engenharia Elétrica & Doutor & DE \\
	Juliana Cristina Braga & Computação Aplicada & Doutora & DE \\
	Karla Vittori & Engenharia Elétrica & Doutora & DE \\
	Luiz Carlos da Silva Rozante & Bioinformática & Doutor & DE \\
	Márcio Katsumi Oikawa & Ciência da Computação & Doutor & DE \\
	Maria das Graças Bruno Marietto & Engenharia Eletrônica & Doutora & DE \\
	Maycon Sambinelli &Ciência da Computação  & Doutor & DE \\
	Monael Pinheiro Ribeiro & Engenharia Eletrônica & Doutor & DE \\
	Nunzio Marco Torrisi & Engenharia Informática & Doutor & DE \\
	Paulo Henrique Pisani & Ciência da Computação e Matemática Computacional & Doutor & DE \\
	Paulo Roberto Miranda Meirelles & Ciência da Computação & Doutor & DE \\
	Raphael Yokoingawa de Camargo & Ciência da Computação & Doutor & DE \\
	Rogério Rossi & Engenharia Elétrica & Doutor & DE \\
	Ronaldo Cristiano Prati & Ciência da Computação e Matemática Computacional & Doutor & DE \\
	Saul de Castro Leite & Modelagem Computacional & Doutor & DE \\
	Sílvia Cristina Dotta & Educação & Doutora & DE \\
	Valério Ramos Batista & Matemática & Doutor & DE \\
	Vera Nagamuta & Ciência da Computação & Doutora & DE \\
	Vinicius Cifú Lopes & Matemática & Doutor & DE \\
	Vladimir Emiliano Moreira Rocha & Engenharia Elétrica & Doutor & DE \\
	Wagner Tanaka Botelho & Engenharia & Doutor & DE \\
	\hline
	
\end{longtable}


\subsection{Núcleo Docente Estruturante}

O Núcleo Docente Estruturante (NDE) do BCC é constituído conforme as orientações da Resolução ConsEPE no 179, de 21 de junho de 2014, que institui o NDE no âmbito dos cursos 
de graduação da UFABC e estabelece suas normas de funcionamento. 

São atribuições do Núcleo Docente Estruturante (NDE).

\begin{enumerate}
	\item [I -] contribuir para a consolidação do perfil profissional do egresso do curso;
	\item [II -] zelar pela integração curricular interdisciplinar entre as diferentes atividades de ensino
	constantes no currículo;
	\item [III -] indicar formas de incentivo ao desenvolvimento de linhas de pesquisa e extensão e sua
	articulação com a pós-graduação, oriundas das necessidades do curso de graduação, das
	exigências do mundo do trabalho, sintonizadas com as políticas públicas próprias à área de
	conhecimento; e
	\item [IV -] zelar pelo cumprimento das Diretrizes Curriculares Nacionais para o Curso e demais
	marcos regulatórios.
\end{enumerate}

A composição atual foi nomeada por meio da Portaria CMCC no. 15/2016, sendo formada 
pelos seguintes docentes:

\begin{itemize}
	\item Márcio Katsumi Oikawa (Presidente)
	\item Carlos da Silva dos Santos
	\item Carlos Alberto Kamienski
	\item Cristiane Maria Sato
	\item Daniel Morgato Martin
	\item Denise Hideko Goya
	\item Diogo Santana Martins
	\item Jerônimo Cordoni Pellegrini
	\item João Marcelo Borovina Josko
\end{itemize}


\newpage
\section {Sistema de Avaliação do Projeto de Curso}

Buscando conhecer, avaliar e aprimorar a qualidade e os compromissos de sua missão, a
UFABC tem implementado mecanismos de avaliação permanente para a efetividade do processo de ensino-aprendizagem, visando compatibilizar a oferta de vagas, os objetivos do curso, o perfil do egresso e a demanda de profissionais no mercado de trabalho para o curso.

Um dos mecanismos adotado pela Coordenação do Curso para avaliação do Projeto
Político Pedagógico do Bacharelado em Ciência da Computação é a análise e o
estabelecimento de ações, a partir dos resultados obtidos pelo Curso e pela Universidade no
Sistema Nacional de Avaliação da Educação Superior (SINAES), regulamentado e instituído
pela Lei n° 10.681, de 14 de abril de 2004.

No Decreto n° 5.773, de 9 de maio de 2006, que dispõe sobre o exercício das funções de
regulação, supervisão e avaliação de Instituições de Educação Superior (IES) e Cursos
superiores de Graduação e Sequenciais no sistema federal de ensino, no seu artigo 1°,
parágrafo 3°, lê-se que a avaliação realizada pelo SINAES constitui referencial básico para os
processos de regulação e supervisão da educação superior, a fim de promover sua qualidade.

No que tange propriamente à estruturação da avaliação estabelecida pelo SINAES, são
considerados três tipos de avaliação:

\begin{enumerate}
	\item Avaliação institucional, que contempla um processo de autoavaliação realizado pela Comissão Própria de Avaliação (CPA) da Instituição de Educação Superior, está já implantada na UFABC, com as seguintes portarias de criação nos últimos anos:
	\begin{enumerate}
		\item Portaria 108, de 28 de fevereiro de 2013, que institui a Comissão Própria de Avaliação e demais portarias correlatas. Disponíveis em
		\url{https://www.ufabc.edu.br/administracao/comissoes/cpa/criacao}. Acesso em: 24 jun. 2022.
		\item Regimento interno da CPA - UFABC. Disponível em \url{https://www.ufabc.edu.br/administracao/comissoes/cpa/regimento-interno}. Acesso em 24 jun. 2022.
	\end{enumerate}
	
	\item Avaliação de curso, que considera um conjunto de avaliações: avaliação dos pares (in loco), avaliação dos estudantes (questionário de Avaliação Discente da Educação Superior – ADES, enviado à amostra selecionada para realização do Exame Nacional de Desempenho de Estudantes - ENADE), avaliação da Coordenação (questionário específico) e dos Professores do Curso e da CPA. Temos os seguintes relatórios produzidos nos últimos anos:
	\begin{enumerate}
		\item Relatório CPA 2022. Disponível em:
		\url{https://www.ufabc.edu.br/images/comissoes/cpa/relatorio_cpa_2022_vfinal_16_04_2022_.pdf}. Acesso em 24 jun. 2022.
		\item Relatório final CPA 2021. Disponível em:
		\url{https://www.ufabc.edu.br/images/comissoes/cpa/relatorio_cpa_2021_final_31_03_2021_entregue.pdf}. Acesso em 24 jun. 2022
		\item Relatório parcial CPA 2021. Disponível em:
		\url{https://www.ufabc.edu.br/images/comissoes/cpa/relatorio_cpa_2020.pdf}. Acesso em: 24 jun. 2022.
		\item Relatório do Grupo de Trabalho sobre Problemas e Oportunidades de Melhoria na Infraestrutura Pedagógica da UFABC. Disponível em \url {https://www.ufabc.edu.br/images/comissoes/cpa/relatorio_gt_infraestrutura_pedagogica.pdf}. Acesso em 24 jun. 2022.
		\item Demais relatórios da CPA - UFABC. Disponíveis em  \url{https://www.ufabc.edu.br/administracao/comissoes/cpa}
	\end{enumerate}
	
	\item Avaliação do Desempenho dos estudantes ingressantes e concluintes, que corresponde à aplicação do ENADE aos estudantes que preenchem os critérios estabelecidos pela legislação vigente (incluem neste exame a prova e os questionários dos alunos, do Coordenador de Curso e da percepção do alunado sobre a prova).
\end{enumerate}

Com o apoio do NDE, os relatórios são utilizados para avaliar a estrutura do curso sob diferentes perspectivas: do discente, do docente, do resultado de exames de acompanhamento externo. Com base nesses elementos, são identificados e discutidos temas levantados sobre pontos positivos e negativos da concepção e execução do curso, como por exemplo:
\begin{itemize}
	\item Adequação da oferta de turmas de disciplinas;
	\item Nível de aproveitamento em disciplinas;
	\item Panorama geral de orientação de alunos para estágios, PGCs, iniciação científica e outras modalidades;
	\item Criação de disciplinas novas com cobertura de assuntos recentes à Computação;
	\item Reformulação de disciplinas;
	\item Adequação de ementas;
	\item Criação de grupos de trabalho;
	\item Outros temas.
\end{itemize}

A aplicação e divulgação dos resultados de discussões realizadas pela coordenação de curso, colegiado de curso e NDE são expostas e deliberadas em reunião plenária, excetuando-se casos em que os temas fogem de seu escopo.


\newpage
\section {Disciplinas}

\subsection{Disciplinas Obrigatórias do BC\&T}
A Tabela \ref{tab:disciplinas_bct}, a seguir, apresenta a lista de todas as disciplinas obrigatórias do BC\&T, que compõe parte do currículo do BCC.

\begin{table}[!h]
	\begin{longtable}{|l|p{.6\textwidth}|c|c|}
		\hline
		Código & Nome da disciplina & Créditos (T-P-I) & Carga-horária\\
		\hline\hline
		BIS0005-15 & Bases Computacionais da Ciência & 2 (0-2-2) & 24h\\
		\hline
		BIJ0207-15 & Bases Conceituais da Energia & 2 (2-0-4) & 24h\\
		\hline
		BIR0004-15 & Bases Epistemológicas da Ciência Moderna & 3 (3-0-4) & 36h\\
		\hline
		BCS0001-15 & Base Experimental das Ciências Naturais & 3 (0-3-2) & 36h\\
		\hline
		BIS0003-15 & Bases Matemáticas & 4 (4-0-5) & 48h\\
		\hline
		BCL0306-15 & Biodiversidade: Interações entre Organismos e Ambiente & 3 (3-0-4) & 36h\\
		\hline
		BCL0308-15 & Bioquímica: Estrutura, Propriedades e Funções de Biomoléculas & 5 (3-2-6) & 60h\\
		\hline
		BIR0603-15 & Ciência, Tecnologia e Sociedade & 3 (3-0-4) & 36h\\
		\hline
		BCM0506-15 & Comunicação e Redes & 3 (3-0-4) & 36h\\
		\hline
		BIK0102-15 & Estrutura da Matéria & 3 (3-0-4) & 36h\\
		\hline
		BIQ0602-15 & Estrutura e Dinâmica Social & 3 (3-0-4) & 36h \\
		\hline
		BIL0304-15 & Evolução e Diversificação da Vida na Terra & 3 (3-0-4) & 36h\\
		\hline
		BCJ0203-15 & Fenômenos Eletromagnéticos & 5 (4-1-6) & 60h\\
		\hline
		BCJ0204-15 & Fenômenos Mecânicos & 5 (4-1-6) & 60h \\
		\hline
		BCJ0205-15 & Fenômenos Térmicos & 4 (3-1-4) & 48h\\
		\hline
		BCK0103-15 & Física Quântica & 3 (3-0-4) & 36h\\
		\hline
		BCN0402-15 & Funções de Uma Variável & 4 (4-0-6) & 48h\\
		\hline
		BCN0407-15 & Funções de Várias Variáveis & 4 (4-0-4) & 48h \\
		\hline
		BCN0404-15 & Geometria Analítica & 3 (3-0-6) & 36h\\
		\hline
		BCK0104-15 & Interações Atômicas e Moleculares & 3 (3-0-4) & 36h\\
		\hline
		BIN0406-15 & Introdução à Probabilidade e Estatística & 3 (3-0-4) & 36h\\
		\hline
		BCN0405-15 & Introdução às Equações Diferenciais Ordinárias & 4 (4-0-4) & 48h \\
		\hline
		BCM0504-15 & Natureza da Informação & 3 (3-0-4) & 36h\\
		\hline
		BCM0505-15 & Processamento da Informação & 5 (3-2-5) & 60h\\
		\hline
		BCS0002-15 & Projeto Dirigido & 2 (0-2-10) & 24h\\
		\hline
		BCL0307-15 & Transformações Químicas & 5 (3-2-6) & 60h\\
		
		\hline
	\end{longtable}
	\caption{Lista de disciplinas obrigatórias do BC\&T}
	\label{tab:disciplinas_bct}
\end{table}


\subsection {Disciplinas Obrigatórias do BCC}

\begin{longtable}{|p{.15\textwidth}|p{.55\textwidth}|p{.15\textwidth}|p{.15\textwidth}|}
	\hline
	Código & Nome da disciplina & Créditos (T-P-I) & Carga-horária\\
	\hline\hline
	MCTB001-17 & Álgebra Linear & 6 (6-0-5) & 72h \\
	\hline
	MCTA001-17 & Algoritmos e Estruturas de Dados I & 4 (2-2-4) & 48h \\
	\hline
	MCTA002-17 & Algoritmos e Estruturas de Dados II & 4 (2-2-4) & 48h \\
	\hline
	MCTA003-17 & Análise de Algoritmos & 4 (4-0-4) & 48h \\
	\hline
	MCTA004-17 & Arquitetura de Computadores & 4 (4-0-4) & 48h \\
	\hline
	MCTA037-17 & Banco de Dados & 4 (3-1-4) & 48h \\
	\hline
	MCTA006-17 & Circuitos Digitais & 4 (3-1-4) & 48h \\
	\hline
	MCTA007-17 & Compiladores & 4 (3-1-4) & 48h \\
	\hline
	MCTA008-17 & Computação Gráfica & 4 (3-1-4) & 48h  \\
	\hline
	MCTA009-13 & Computadores, Ética e Sociedade & 2 (2-0-4) & 24h \\
	\hline
	MCTA033-15 & Engenharia de Software & 4 (4-0-4) & 48h \\
	\hline
	MCTA014-15 & Inteligência Artificial & 4 (3-1-4) & 48h \\
	\hline
	MCTA015-13 & Linguagens Formais e Automata & 4 (3-1-4) & 48h \\
	\hline
	NHI2049-13 & Lógica Básica & 4 (4-0-4) & 48h \\
	\hline
	MCTB019-17 & Matemática Discreta & 4 (4-0-4) & 48h \\
	\hline
	MCTA016-13 & Paradigmas de Programação & 4 (2-2-4) & 48h \\
	\hline
	MCTA028-15 & Programação Estruturada & 4 (2-2-4) & 48h \\
	\hline
	MCTA017-17 & Programação Matemática & 4 (3-1-4) & 48h \\
	\hline
	MCTA018-13 & Programação Orientada a Objetos & 4 (2-2-4) & 48h \\
	\hline
	MCTA029-17 & Projeto de Graduação em Computação I & 8 (0-8-8) & 96h \\
	\hline
	MCTA030-17 & Projeto de Graduação em Computação II & 8 (0-8-8) & 96h \\
	\hline
	MCTA031-17 & Projeto de Graduação em Computação III & 8 (0-8-8) & 96h \\
	\hline
	MCTA022-17 & Redes de Computadores & 4 (3-1-4) & 48h \\
	\hline
	MCTA023-17 & Segurança de Dados & 4 (3-1-4) & 48h \\
	\hline
	MCTA024-13 & Sistemas Digitais & 4 (2-2-4) & 48h \\
	\hline
	MCTA025-13 & Sistemas Distribuídos & 4 (3-1-4) & 48h \\
	\hline
	MCTA026-13 & Sistemas Operacionais & 4 (3-1-4) & 48h\\
	\hline
	MCTA027-17 & Teoria dos Grafos & 4 (3-1-4) & 48h\\
	\hline
\end{longtable}


\newpage
\section{Disciplinas de opção limitada}

Lista de disciplinas de opção limitada do Bacharelado em Ciência da Computação
atualizado em julho de 2022.

\begin{longtable}{|p{.1\textwidth}|p{.6\textwidth}|p{.15\textwidth}|p{.15\textwidth}|}
    \hline
    Código & Nome da disciplina & Créditos (TPI) & Carga-horária\\
    \hline\hline
\end{longtable}


\newpage
\section{Regras de transição}
\label{sec:regras_transicao}

A matriz curricular de 2023 entrará em vigor assim que aprovada por todos os
órgãos deliberativos de acordo com a Resolução ConsEPE nº 140 e será plenamente
oferecida para os ingressantes a partir do ano de 2023 na Universidade Federal
do ABC.

Os alunos que ingressaram na UFABC anteriormente a 2023 poderão integralizar o
curso de acordo com a matriz curricular de 2023 ou de acordo com a matriz
curricular do projeto pedagógico vigente no seu ano de ingresso.
Caberá ao aluno realizar a análise da sua situação com relação ao coeficiente
de progressão e decidir por qual matriz pretende obter o grau de bacharel em
Ciência da Computação.
Será necessário cumprir em sua totalidade uma das matrizes para ter direito à
integralização do curso.

Os ingressantes anteriores a 2023 que optem pela integralização de acordo com
o PPC 2023 devem simplesmente cursar as disciplinas obrigatórias constantes na
matriz de 2023.
Todas as disciplinas já cursadas na matriz anterior serão convalidadas
automaticamente, com exceção das seguintes:
\begin{itemize}
    \item A disciplina ``NHI2049-13 -- Lógica Básica'' será convalidada como de
    opção livre;
    \item A disciplina ``MCTA027-17 -- Teoria dos Grafos'' será convalidada
    como ``MCTA027-23 -- Algoritmos em Grafos'';
    \item As disciplinas de ``Estágio Supervisionado em Computação'' serão
    convalidadas como de opção limitada. Essas disciplinas deixaram de ser
    ofertadas em 2019 e estão extintas desde a matriz de 2017.
\end{itemize}

Os ingressantes anteriores a 2023 que optem pela integralização de acordo com o
seu PPC de ingresso também devem simplesmente continuar cursando as disciplinas
que serão ofertadas a partir de 2023, uma vez que todas as disciplinas
obrigatórias dos PPCs anteriores permanecem obrigatórias no PPC de 2023, com as
seguintes observações:
\begin{itemize}
    \item A disciplina ``MCTA016-23 -- Programação Funcional'' será convalidada
    como ``MCTA016-13 -- Paradigmas de Programação'';
    \item A disciplina ``MCTA027-23 -- Algoritmos em Grafos'' será convalidada
    como ``MCTA027-17 -- Teoria dos Grafos'';
    \item A disciplina ``MCTAXXX-23 -- Matemática Discreta II'' será
    convalidada como ``NHI2049-13 -- Lógica Básica''.  Alternativamente, a
    disciplina ``NHI2049-13 -- Lógica Básica'' pode ser cursada de acordo com
    oferta pelo curso de Filosofia;
    \item A disciplina ``MCTAXXX-23 -- Metodologia e Escrita Científica para
    Computação'' será convalidada como opção livre/limitada.
\end{itemize}
Com relação às disciplinas de opção limitada, a convalidação é direta, com
observação para duas disciplinas que tiveram o nome modificado:
\begin{itemize}
    \item A disciplina ``MCZA036-17 -- Análise de Algoritmos II'' passou a se
    chamar ``MCZA036-23 -- Análise de Algoritmos III'';
    \item A disciplina ``MCZA014-17 -- Métodos de Otimização'' passou a se
    chamar ``MCZA014-23 -- Otimização Não-linear''.
\end{itemize}

Detalhes das convalidações podem ser encontrados nas
Tabelas~\ref{tab:convalidacoes_obrigatorias}
e~\ref{tab:convalidacoes_limitadas}.

Os casos omissos serão resolvidos pela coordenação do curso, representada pelo
seu coordenador, com o apoio da Pró-Reitoria de Graduação.

\begin{table}[h!]
\caption{Tabela de convalidações das disciplinas obrigatórias.}
\label{tab:convalidacoes_obrigatorias}

\centering
{\footnotesize
\begin{tabular}{|c|p{.25\textwidth}|c||c|p{.25\textwidth}|c|}
    \hline
    \multicolumn{3}{|c||}{\bf 2017} & \multicolumn{3}{|c|}{\bf 2023} \\ 
    \hline
    
    \textbf{Sigla} & \textbf{Disciplina} & \textbf{Créditos} & \textbf{Sigla} & \textbf{Disciplina} & \textbf{Créditos} \\
    \hline\hline
    
    MCTB001-17 & Álgebra Linear & 6 & MCTB001-23 & Álgebra Linear & 6 \\ \hline
    MCTA001-17 & Algoritmos e Estruturas de Dados I & 4 & MCTA001-23 & Algoritmos e Estruturas de Dados I & 4 \\ \hline
    MCTA002-17 & Algoritmos e Estruturas de Dados II & 4 & MCTA002-23 & Algoritmos e Estruturas de Dados II & 4 \\ \hline
    MCTA003-17 & Análise de Algoritmos & 4 & MCTA003-23 & Análise de Algoritmos I & 4 \\
    & & & MCTAXXX-23 & Análise de Algoritmos II & 4 \\ \hline
    MCTA004-17 & Arquitetura de Computadores & 4 & MCTA004-23 & Arquitetura de Computadores & 4 \\ \hline
    MCTA037-17 & Banco de Dados & 4 & MCTA037-23 & Modelagem de Banco de Dados & 4 \\ \hline
    MCTA006-17 & Circuitos Digitais & 4 & MCTA006-23 & Circuitos Digitais & 4 \\ \hline
    MCTA007-17 & Compiladores & 4 & MCTA007-23 & Compiladores e Interpretadores & 4 \\ \hline
    MCTA008-17 & Computação Gráfica & 4 & MCTA008-23 & Computação Gráfica & 4 \\ \hline
    MCTA009-13 & Computadores, Ética e Sociedade & 2 & MCTA009-23 & Computadores, Ética e Sociedade & 2 \\ \hline
    MCTA033-15 & Engenharia de Software & 4 & MCTA033-23 & Engenharia de Software & 4 \\ \hline
    MCTA014-15 & Inteligência Artificial & 4 & MCTA014-23 & Inteligência Artificial & 4 \\ \hline
    MCTA015-13 & Linguagens Formais e Automata & 4 & MCTA015-23 & Linguagens Formais e Autômatos & 4 \\ \hline
    NHI2049-13 & Lógica Básica & 4 & NHI2049-XX & Lógica Básica & 4 \\
    & & & MCTAXXX-23 & ou Matemática Discreta II & 4 \\ \hline
    MCTB019-17 & Matemática Discreta & 4 & MCTB019-23 & Matemática Discreta & 4 \\ \hline
    MCTA016-13 & Paradigmas de Programação & 4 & MCTA016-23 & Programação Funcional & 4 \\ \hline
    MCTA028-15 & Programação Estruturada & 4 & MCTA028-23 & Programação Estruturada & 4 \\ \hline
    MCTA017-17 & Programação Matemática & 4 & MCTA017-23 & Otimização Linear & 4 \\ \hline
    MCTA018-13 & Programação Orientada a Objetos & 4 & MCTA018-23 & Programação Orientada a Objetos & 4 \\ \hline
    MCTA029-17 & Projeto de Graduação em Computação I & 8 & MCTA029-23 & Trabalho de Conclusão de Curso I & 4 \\ \hline
    MCTA030-17 & Projeto de Graduação em Computação II & 8 & MCTA030-23 & Trabalho de Conclusão de Curso II & 4 \\ \hline
    MCTA031-17 & Projeto de Graduação em Computação III & 8 & MCTA031-23 & Trabalho de Conclusão de Curso III & 4 \\ \hline
    MCTA022-17 & Redes de Computadores & 4 & MCTA022-23 & Redes de Computadores & 4 \\ \hline
    MCTA023-17 & Segurança de Dados & 4 & MCTA023-23 & Segurança de Dados & 4 \\ \hline
    MCTA024-13 & Sistemas Digitais & 4 & MCTA024-23 & Sistemas Digitais & 4 \\ \hline
    MCTA025-13 & Sistemas Distribuídos & 4 & MCTA025-23 & Sistemas Distribuídos & 4 \\ \hline
    MCTA026-13 & Sistemas Operacionais & 4 & MCTA026-23 & Sistemas Operacionais & 4 \\ \hline
    MCTA027-17 & Teoria dos Grafos & 4 & MCTA027-23 & Algoritmos em Grafos & 4 \\ \hline
     & Não há & & MCTAXXX-23 & Metodologia e escrita científica para computação & 2 \\ \hline
\end{tabular}
}
\end{table}

{\footnotesize
\begin{longtable}{|c|p{.25\textwidth}|c||c|p{.25\textwidth}|c|}
\caption{Tabela de convalidações das disciplinas de opção limitada.}
\label{tab:convalidacoes_limitadas} \\

\hline
\multicolumn{3}{|c||}{\bf 2017} & \multicolumn{3}{|c|}{\bf 2023} \\ 
\hline

\textbf{Sigla} & \textbf{Disciplina} & \textbf{Créditos} & \textbf{Sigla} & \textbf{Disciplina} & \textbf{Créditos} \\
\hline\hline

MCZA035-17 & Algoritmos Probabilísticos & 4 & MCZA035-17 & Algoritmos Probabilísticos & 4\\ \hline
MCZA036-17 & Análise de Algoritmos II & 4 & MCZA036-23 & Análise de Algoritmos III & 4\\ \hline
MCZA001-13 & Análise de Projetos & 2 & MCZA001-13 & Análise de Projetos & 2\\ \hline
MCTB007-17 & Anéis e Corpos & 4 & MCTB007-17 & Anéis e Corpos & 4\\ \hline
MCZA002-17 & Aprendizado de Máquina & 4 & MCZA002-17 & Aprendizado de Máquina & 4\\ \hline
MCZA003-17 & Arquitetura de Computadores de Alto Desempenho & 4 & MCZA003-17 & Arquitetura de Computadores de Alto Desempenho & 4\\ \hline
MCZA004-13 & Avaliação de Desempenho de Redes & 4 & MCZA004-13 & Avaliação de Desempenho de Redes & 4\\ \hline
MCZA005-17 & Banco de Dados de Apoio à Tomada de Decisão & 4 & MCZA005-17 & Banco de Dados de Apoio à Tomada de Decisão & 4\\ \hline
MCTB009-17 & Cálculo Numérico & 4 & MCTB009-17 & Cálculo Numérico & 4\\ \hline
MCZA037-17 & Combinatória Extremal & 4 & MCZA037-17 & Combinatória Extremal & 4\\ \hline
MCZA006-17 & Computação Evolutiva e Conexionista & 4 & MCZA006-17 & Computação Evolutiva e Conexionista & 4\\ \hline
ESZG013-17 & Empreendedorismo & 4 & ESZG013-17 & Empreendedorismo & 4\\ \hline
MCZA007-13 & Empreendedorismo e Desenvolvimento de Negócios & 4 & MCZA007-13 & Empreendedorismo e Desenvolvimento de Negócios & 4\\ \hline
MCZA051-17 & Estágio Supervisionado em Computação & 4 & MCZA051-17 & Estágio Supervisionado em Computação & 4\\ \hline
ESZI030-17 & Gerenciamento e Interoperabilidade de Redes & 4 & ESZI030-17 & Gerenciamento e Interoperabilidade de Redes & 4\\ \hline
MCZA016-17 & Gestão de projetos de software & 4 & MCZA016-17 & Gestão de projetos de software & 4\\ \hline
ESZG019-17 & Gestão Estratégica e Organizacional & 2 & ESZG019-17 & Gestão Estratégica e Organizacional & 2\\ \hline
MCTB018-17 & Grupos & 4 & MCTB018-17 & Grupos & 4\\ \hline
MCZB012-13 & Inferência Estatística & 4 & MCZB012-13 & Inferência Estatística & 4\\ \hline
ESZI013-17 & Informática Industrial & 4 & ESZI013-17 & Informática Industrial & 4\\ \hline
MCZA008-17 & Interação Humano-Computador & 4 & MCZA008-17 & Interação Humano-Computador & 4\\ \hline
ESZB022-17 & Introdução à Bioinformática & 4 & ESZB022-17 & Introdução à Bioinformática & 4\\ \hline
MCZB015-13 & Introdução à Criptografia & 4 & MCZB015-13 & Introdução à Criptografia & 4\\ \hline
MCZB018-13 & Introdução à Modelagem e Processos Estocásticos & 4 & MCZB018-13 & Introdução à Modelagem e Processos Estocásticos & 4\\ \hline
MCTC021-15 & Introdução à Neurociência Computacional & 4 & MCTC021-15 & Introdução à Neurociência Computacional & 4\\ \hline
MCZA032-17 & Introdução à Programação de Jogos & 4 & MCZA032-17 & Introdução à Programação de Jogos & 4\\ \hline
ESZI034-17 & Jogos Digitais: Aspectos Técnicos e Aplicações & 4 & ESZI034-17 & Jogos Digitais: Aspectos Técnicos e Aplicações & 4\\ \hline
MCZA010-13 & Laboratório de Engenharia de Software & 4 & MCZA010-13 & Laboratório de Engenharia de Software & 4\\ \hline
MCZA011-17 & Laboratório de Redes & 4 & MCZA011-17 & Laboratório de Redes & 4\\ \hline
MCZA012-13 & Laboratório de Sistemas Operacionais & 4 & MCZA012-13 & Laboratório de Sistemas Operacionais & 4\\ \hline
MCZA013-13 & Lógicas não Clássicas & 4 & MCZA013-13 & Lógicas não Clássicas & 4\\ \hline
MCZA014-17 & Métodos de Otimização & 4 & MCZA014-23 & Otimização não-linear & 4\\ \hline
MCZA015-13 & Mineração de Dados & 4 & MCZA015-13 & Mineração de Dados & 4\\ \hline
MCZA052-22 & Vizualização de Dados e Informações & 4 & MCZA052-22 & Vizualização de Dados e Informações & 4\\ \hline
ESTG013-17 & Pesquisa Operacional & 4 & ESTG013-17 & Pesquisa Operacional & 4\\ \hline
ESZI022-17 & Planejamento de Redes de Informação & 4 & ESZI022-17 & Planejamento de Redes de Informação & 4\\ \hline
MCZA038-17 & Prática Avançada de Programação A & 4 & MCZA038-17 & Prática Avançada de Programação A & 4\\ \hline
MCZA038-17 & Prática Avançada de Programação B & 4 & MCZA038-17 & Prática Avançada de Programação B & 4\\ \hline
MCZA038-17 & Prática Avançada de Programação C & 4 & MCZA038-17 & Prática Avançada de Programação C & 4\\ \hline
MCZA041-17 & Processamento de Imagens Utilizando GPU & 4 & MCZA041-17 & Processamento de Imagens Utilizando GPU & 4\\ \hline
MCZA017-13 & Processamento de Linguagem Natural & 4 & MCZA017-13 & Processamento de Linguagem Natural & 4\\ \hline
MCTC022-15 & Processamento de Sinais Neurais & 4 & MCTC022-15 & Processamento de Sinais Neurais & 4\\ \hline
MCZA018-17 & Processamento Digital de Imagens & 4 & MCZA018-17 & Processamento Digital de Imagens & 4\\ \hline
MCZA042-17 & Processo e Desenvolvimento de Softwares Educacionais & 4 & MCZA042-17 & Processo e Desenvolvimento de Softwares Educacionais & 4\\ \hline
MCZA033-17 & Programação Avançada para Dispositivos Móveis & 4 & MCZA033-17 & Programação Avançada para Dispositivos Móveis & 4\\ \hline
ESZI033-17 & Programação de Dispositivos Móveis & 2 & ESZI033-17 & Programação de Dispositivos Móveis & 2\\ \hline
MCZA019-17 & Programação para Web & 4 & MCZA019-17 & Programação para Web & 4\\ \hline
MCZA020-13 & Programação Paralela & 4 & MCZA020-13 & Programação Paralela & 4\\ \hline
MCZA034-17 & Programação Segura & 4 & MCZA034-17 & Programação Segura & 4\\ \hline
MCZA021-17 & Projeto de Redes & 4 & MCZA021-17 & Projeto de Redes & 4\\ \hline
MCZA022-17 & Projeto Interdisciplinar & 4 & MCZA022-17 & Projeto Interdisciplinar & 4\\ \hline
MCZA023-17 & Redes Convergentes & 4 & MCZA023-17 & Redes Convergentes & 4\\ \hline
ESZI029-17 & Redes WAN de Banda Larga & 4 & ESZI029-17 & Redes WAN de Banda Larga & 4\\ \hline
MCZA024-17 & Redes sem Fio & 4 & MCZA024-17 & Redes sem Fio & 4\\ \hline
MCZA044-17 & Robótica e Sistemas Inteligentes & 4 & MCZA044-17 & Robótica e Sistemas Inteligentes & 4\\ \hline
MCZA045-17 & Robótica Educacional & 4 & MCZA045-17 & Robótica Educacional & 4\\ \hline
MCZA025-13 & Segurança em Redes & 4 & MCZA025-13 & Segurança em Redes & 4\\ \hline
MCZA046-17 & Semântica de Linguagem de Programação & 4 & MCZA046-17 & Semântica de Linguagem de Programação & 4\\ \hline
MCZA026-17 & Sistemas de Gerenciamento de Banco de Dados & 4 & MCZA026-17 & Sistemas de Gerenciamento de Banco de Dados & 4\\ \hline
MCZA027-17 & Sistemas de Informação & 4 & MCZA027-17 & Sistemas de Informação & 4\\ \hline
ESZI014-17 & Sistemas Inteligentes & 4 & ESZI014-17 & Sistemas Inteligentes & 4\\ \hline
MCZA028-13 & Sistemas Multiagentes & 4 & MCZA028-13 & Sistemas Multiagentes & 4\\ \hline
MCZA029-13 & Sistemas Multimídia & 4 & MCZA029-13 & Sistemas Multimídia & 4\\ \hline
MCZA047-17 & Sistemas Multi-Robôs Sociais & 4 & MCZA047-17 & Sistemas Multi-Robôs Sociais & 4\\ \hline
MCZA050-15 & Técnicas Avançadas de Programação & 4 & MCZA050-15 & Técnicas Avançadas de Programação & 4\\ \hline
NHZ5019-15 & Tecnologias da Informação e Comunicação na Educação & 4 & NHZ5019-15 & Tecnologias da Informação e Comunicação na Educação & 4\\ \hline
MCZB033-17 & Teoria da Recursão e Computabilidade & 4 & MCZB033-17 & Teoria da Recursão e Computabilidade & 4\\ \hline
MCZA048-17 & Teoria Espectral de Grafos & 4 & MCZA048-17 & Teoria Espectral de Grafos & 4\\ \hline
MCZA049-17 & Tópicos Emergentes em Banco de Dados & 4 & MCZA049-17 & Tópicos Emergentes em Banco de Dados & 4\\ \hline
MCZA030-17 & Vida Artificial na Computação & 4 & MCZA030-17 & Vida Artificial na Computação & 4\\ \hline
ESZA019-17 & Visão Computacional & 4 & ESZA019-17 & Visão Computacional & 4\\ \hline
MCZA031-13 & Web Semântica & 4 & MCZA031-13 & Web Semântica & 4\\ \hline

\end{longtable}
}



\end{document}


